\worldchapter{\trans{items_title}}{appendix-items}
\label{ch:appendix-items}

\begin{multicols}{2}


    \section{Erste Hilfe}



    \subsection*{Bandagen}

    Ermöglicht die Anwendung von Erster Hilfe.

    \begin{tabular}{@{}ll@{}}
        \textbf{\trans{weight_label}:} & 0.30 \\
        \textbf{\trans{price_label}:} & 5 \\
        \textbf{\trans{rarity_label}:} & Gewöhnlich \\
        \textbf{\trans{concealment_label}:} & 0 \\
    \end{tabular}

    \wbif{phasesix}{
        \vspace{3mm}
\faIcon{cube}\hspace{0.5em}    }{}


    \subsection*{Notfallkoffer}

    Ermöglicht die Anwendung von Erster Hilfe, und heilt zusätzlich 1W3 weitere Wunden.

Enthält 5 Anwendungen.

    \begin{tabular}{@{}ll@{}}
        \textbf{\trans{weight_label}:} & 5.00 \\
        \textbf{\trans{price_label}:} & 100 \\
        \textbf{\trans{rarity_label}:} & Gewöhnlich \\
        \textbf{\trans{concealment_label}:} & 2 \\
        \textbf{\trans{charges_label}:} & 5\\
    \end{tabular}

    \wbif{phasesix}{
        \vspace{3mm}
\faIcon{umbrella}\hspace{0.5em}\faIcon{flag}\hspace{0.5em}\faIcon{plane}\hspace{0.5em}\faIcon{microchip}\hspace{0.5em}\faIcon{space-shuttle}\hspace{0.5em}    }{}


    \subsection*{Adrenalinspritze}

    Bei der Anwendung erzeugt die Adrenalinspritze 1W6+2 Boost.

    \begin{tabular}{@{}ll@{}}
        \textbf{\trans{weight_label}:} & 0.20 \\
        \textbf{\trans{price_label}:} & 200 \\
        \textbf{\trans{rarity_label}:} & Gewöhnlich \\
        \textbf{\trans{concealment_label}:} & 0 \\
    \end{tabular}

    \wbif{phasesix}{
        \vspace{3mm}
\faIcon{flag}\hspace{0.5em}\faIcon{plane}\hspace{0.5em}\faIcon{microchip}\hspace{0.5em}\faIcon{space-shuttle}\hspace{0.5em}    }{}


    \section{Tränke und Gifte}



    \subsection*{Einfache Wundtinktur}

    Bei erfolgreicher Anwendung mit \textit{Erster Hilfe} und einer Bandage heilt die Bandage 1W3 Wunden zusätzlich.

    \begin{tabular}{@{}ll@{}}
        \textbf{\trans{weight_label}:} & 0.80 \\
        \textbf{\trans{price_label}:} & 30 \\
        \textbf{\trans{rarity_label}:} & Gewöhnlich \\
        \textbf{\trans{concealment_label}:} & 0 \\
    \end{tabular}

    \wbif{phasesix}{
        \vspace{3mm}
\faIcon{cube}\hspace{0.5em}    }{}


    \subsection*{Morphin}

    Eine Flasche Morphinflüssigkeit

    \begin{tabular}{@{}ll@{}}
        \textbf{\trans{weight_label}:} & 0.01 \\
        \textbf{\trans{price_label}:} & 20 \\
        \textbf{\trans{rarity_label}:} & Selten \\
        \textbf{\trans{concealment_label}:} & 1 \\
    \end{tabular}

    \wbif{phasesix}{
        \vspace{3mm}
\faIcon{umbrella}\hspace{0.5em}\faIcon{flag}\hspace{0.5em}\faIcon{plane}\hspace{0.5em}\faIcon{microchip}\hspace{0.5em}    }{}


    \subsection*{grobschlächtige Spritze}

    gefüllt mit einer seltsamen Flüssigkeit

    \begin{tabular}{@{}ll@{}}
        \textbf{\trans{weight_label}:} & 1.00 \\
        \textbf{\trans{price_label}:} & 10 \\
        \textbf{\trans{rarity_label}:} & Einmalig \\
        \textbf{\trans{concealment_label}:} & 0 \\
    \end{tabular}

    \wbif{phasesix}{
        \vspace{3mm}
\faIcon{umbrella}\hspace{0.5em}\faIcon{industry}\hspace{0.5em}\faIcon{ankh}\hspace{0.5em}\faIcon{syringe}\hspace{0.5em}\faIcon{spider}\hspace{0.5em}    }{}


    \section{Werfbares}



    \subsection*{Blendgranate}

    Kann bis zu 30 Meter weit geworfen werden.
Entfernt im Radius von W6 Metern alle noch vorhandenen Aktionen von Teilnehmern am Kampf. Getroffene Teilnehmer erhalten am Anfang ihrer nächsten Kampfrunde nur die Hälfte ihrer Aktionen (abgerundet).

    \begin{tabular}{@{}ll@{}}
        \textbf{\trans{weight_label}:} & 0.40 \\
        \textbf{\trans{price_label}:} & 100 \\
        \textbf{\trans{rarity_label}:} & Gewöhnlich \\
        \textbf{\trans{concealment_label}:} & 1 \\
    \end{tabular}

    \wbif{phasesix}{
        \vspace{3mm}
\faIcon{flag}\hspace{0.5em}\faIcon{plane}\hspace{0.5em}\faIcon{microchip}\hspace{0.5em}\faIcon{space-shuttle}\hspace{0.5em}    }{}


    \subsection*{Explosivgranate}

    Kann bis zu 30 Meter weit geworfen werden.
Verursacht im Radius von W6 Metern an jedem Teilnehmer am Kampf 3 Treffer mit je zwei Wunden und Durchschlag 1.

    \begin{tabular}{@{}ll@{}}
        \textbf{\trans{weight_label}:} & 0.50 \\
        \textbf{\trans{price_label}:} & 120 \\
        \textbf{\trans{rarity_label}:} & Gewöhnlich \\
        \textbf{\trans{concealment_label}:} & 1 \\
    \end{tabular}

    \wbif{phasesix}{
        \vspace{3mm}
\faIcon{flag}\hspace{0.5em}\faIcon{plane}\hspace{0.5em}\faIcon{microchip}\hspace{0.5em}\faIcon{space-shuttle}\hspace{0.5em}    }{}


    \subsection*{Tränengasgranate}

    Kann bis zu 30 Meter weit geworfen werden.
Wenn die Tränengasgranate geworfen wird entwickelt sich am Ende der Kampfrunde eine Gaswolke von W6 Metern Durchmesser. Innerhalb der Gaswolke verringern sich die Aktionen jedes Charakters auf 1 pro Runde, Wahrnehmung und Willenskraft sind um 3 verringert.

Die Gaswolke bleibt 1W6+2 Runden bestehen.

    \begin{tabular}{@{}ll@{}}
        \textbf{\trans{weight_label}:} & 0.20 \\
        \textbf{\trans{price_label}:} & 200 \\
        \textbf{\trans{rarity_label}:} & Außergewöhnlich \\
        \textbf{\trans{concealment_label}:} & 1 \\
    \end{tabular}

    \wbif{phasesix}{
        \vspace{3mm}
\faIcon{flag}\hspace{0.5em}\faIcon{plane}\hspace{0.5em}\faIcon{microchip}\hspace{0.5em}    }{}


    \subsection*{Giftgasgranate}

    Kann bis zu 30 Meter weit geworfen werden.
Wenn die Giftgasgranate geworfen wird entwickelt sich am Ende der Kampfrunde eine Gaswolke von W6 Metern Durchmesser. Innerhalb der Gaswolke verringert sich die Anzahl der Aktionen um 1, und der Charakter erhält den Statuswert Vergiftet 2.

Die Gaswolke bleibt 1W6+2 Runden bestehen.

    \begin{tabular}{@{}ll@{}}
        \textbf{\trans{weight_label}:} & 0.20 \\
        \textbf{\trans{price_label}:} & 350 \\
        \textbf{\trans{rarity_label}:} & Außergewöhnlich \\
        \textbf{\trans{concealment_label}:} & 1 \\
    \end{tabular}

    \wbif{phasesix}{
        \vspace{3mm}
\faIcon{flag}\hspace{0.5em}\faIcon{plane}\hspace{0.5em}\faIcon{microchip}\hspace{0.5em}    }{}


    \subsection*{Rauchgranate}

    Kann bis zu 30 Meter weit geworfen werden.
Wenn die Rauchgranate geworfen wird entwickelt sich am Ende der Kampfrunde eine Rauchwolke von 2W6 Metern Radius. Die Rauchwolke behindert jegliche Sicht und Wahrnehmung.

Die Rauchwolke bleibt für 1W6+2 Runden bestehen.

    \begin{tabular}{@{}ll@{}}
        \textbf{\trans{weight_label}:} & 0.20 \\
        \textbf{\trans{price_label}:} & 50 \\
        \textbf{\trans{rarity_label}:} & Gewöhnlich \\
        \textbf{\trans{concealment_label}:} & 1 \\
    \end{tabular}

    \wbif{phasesix}{
        \vspace{3mm}
\faIcon{flag}\hspace{0.5em}\faIcon{plane}\hspace{0.5em}\faIcon{microchip}\hspace{0.5em}    }{}


    \subsection*{Molotov Cocktail}

    Kann bis zu 30 Meter weit geworfen werden.
Verursacht im Radius von W6 Metern an jedem Teilnehmer am Kampf einen Treffer mit einer Wunde und Durchschlag 0. Verursacht bei jedem Getroffenen Brennend 2. Entzündet brennbares.

Das Feuer bleibt für 1W6+2 Runden bestehen.

    \begin{tabular}{@{}ll@{}}
        \textbf{\trans{weight_label}:} & 1.00 \\
        \textbf{\trans{price_label}:} & 20 \\
        \textbf{\trans{rarity_label}:} & Gewöhnlich \\
        \textbf{\trans{concealment_label}:} & 1 \\
    \end{tabular}

    \wbif{phasesix}{
        \vspace{3mm}
\faIcon{umbrella}\hspace{0.5em}\faIcon{flag}\hspace{0.5em}\faIcon{plane}\hspace{0.5em}\faIcon{microchip}\hspace{0.5em}\faIcon{space-shuttle}\hspace{0.5em}    }{}


    \subsection*{Wurfnetz}

    Das Wurfnetz kann im Kampf geworfen werden, um den Gegner im Netz zu fangen. 

Ist der Werfen-Wurf gelungen, so gilt der Gegner als gefangen. Er benötigt einen Geschick-Wurf, um sich aus dem Netz zu befreien (1 Aktion). Solange der Gegner im Netz gefangen ist, kann er sich nicht bewegen, alle Aktionen sind schwere Proben.

    \begin{tabular}{@{}ll@{}}
        \textbf{\trans{weight_label}:} & 1.00 \\
        \textbf{\trans{price_label}:} & 30 \\
        \textbf{\trans{rarity_label}:} & Gewöhnlich \\
        \textbf{\trans{concealment_label}:} & 0 \\
    \end{tabular}

    \wbif{phasesix}{
        \vspace{3mm}
\faIcon{cube}\hspace{0.5em}    }{}


    \section{Gefäße}



    \subsection*{Keramikflasche}


    \begin{tabular}{@{}ll@{}}
        \textbf{\trans{weight_label}:} & 0.20 \\
        \textbf{\trans{price_label}:} & 10 \\
        \textbf{\trans{rarity_label}:} & Gewöhnlich \\
        \textbf{\trans{concealment_label}:} & 0 \\
    \end{tabular}

    \wbif{phasesix}{
        \vspace{3mm}
\faIcon{cube}\hspace{0.5em}    }{}


    \subsection*{Phiole}

    Eine Phiole aus Glas

    \begin{tabular}{@{}ll@{}}
        \textbf{\trans{weight_label}:} & 0.10 \\
        \textbf{\trans{price_label}:} & 20 \\
        \textbf{\trans{rarity_label}:} & Gewöhnlich \\
        \textbf{\trans{concealment_label}:} & 0 \\
    \end{tabular}

    \wbif{phasesix}{
        \vspace{3mm}
\faIcon{cube}\hspace{0.5em}    }{}


    \subsection*{Ledertasche}


    \begin{tabular}{@{}ll@{}}
        \textbf{\trans{weight_label}:} & 0.80 \\
        \textbf{\trans{price_label}:} & 15 \\
        \textbf{\trans{rarity_label}:} & Gewöhnlich \\
        \textbf{\trans{concealment_label}:} & 2 \\
    \end{tabular}

    \wbif{phasesix}{
        \vspace{3mm}
\faIcon{cube}\hspace{0.5em}    }{}


    \subsection*{Tuchbeutel}

    Der Tuchbeutel kann genutzt werden um Gegenstände darin zu lagern oder zu transportieren.

    \begin{tabular}{@{}ll@{}}
        \textbf{\trans{weight_label}:} & 0.50 \\
        \textbf{\trans{price_label}:} & 5 \\
        \textbf{\trans{rarity_label}:} & Gewöhnlich \\
        \textbf{\trans{concealment_label}:} & 0 \\
    \end{tabular}

    \wbif{phasesix}{
        \vspace{3mm}
\faIcon{cube}\hspace{0.5em}    }{}


    \subsection*{Lederranzen}

    Ein bequem zu tragender Lederranzen, in dem Gegenstände verstaut werden können.

    \begin{tabular}{@{}ll@{}}
        \textbf{\trans{weight_label}:} & 2.00 \\
        \textbf{\trans{price_label}:} & 20 \\
        \textbf{\trans{rarity_label}:} & Gewöhnlich \\
        \textbf{\trans{concealment_label}:} & 1 \\
    \end{tabular}

    \wbif{phasesix}{
        \vspace{3mm}
\faIcon{cube}\hspace{0.5em}    }{}


    \subsection*{Geldbeutel}

    Ein Geldbeutel, gerade groß genug um dein Vermögen zu tragen.

    \begin{tabular}{@{}ll@{}}
        \textbf{\trans{weight_label}:} & 0.20 \\
        \textbf{\trans{price_label}:} & 10 \\
        \textbf{\trans{rarity_label}:} & Gewöhnlich \\
        \textbf{\trans{concealment_label}:} & 0 \\
    \end{tabular}

    \wbif{phasesix}{
        \vspace{3mm}
\faIcon{umbrella}\hspace{0.5em}\faIcon{flag}\hspace{0.5em}\faIcon{plane}\hspace{0.5em}\faIcon{microchip}\hspace{0.5em}    }{}


    \subsection*{Pergamenthülle}

    Hierin sind deine Dokumente sicher! Eine lederne, wasserfeste Hülle um Pergamente oder Dokumente zu verstauen.

    \begin{tabular}{@{}ll@{}}
        \textbf{\trans{weight_label}:} & 0.20 \\
        \textbf{\trans{price_label}:} & 40 \\
        \textbf{\trans{rarity_label}:} & Gewöhnlich \\
        \textbf{\trans{concealment_label}:} & 0 \\
    \end{tabular}

    \wbif{phasesix}{
        \vspace{3mm}
\faIcon{cube}\hspace{0.5em}    }{}


    \subsection*{Sack}

    Ein Sack aus Leinen, groß genug um viele Gegenstände zu transportieren.

    \begin{tabular}{@{}ll@{}}
        \textbf{\trans{weight_label}:} & 1.00 \\
        \textbf{\trans{price_label}:} & 10 \\
        \textbf{\trans{rarity_label}:} & Gewöhnlich \\
        \textbf{\trans{concealment_label}:} & 1 \\
    \end{tabular}

    \wbif{phasesix}{
        \vspace{3mm}
\faIcon{cube}\hspace{0.5em}    }{}


    \subsection*{Korb}

    In diesem Korb kann man Gegenstände oder anderen Objekte transportieren.

    \begin{tabular}{@{}ll@{}}
        \textbf{\trans{weight_label}:} & 1.00 \\
        \textbf{\trans{price_label}:} & 10 \\
        \textbf{\trans{rarity_label}:} & Gewöhnlich \\
        \textbf{\trans{concealment_label}:} & 0 \\
    \end{tabular}

    \wbif{phasesix}{
        \vspace{3mm}
\faIcon{cube}\hspace{0.5em}    }{}


    \subsection*{Packsattel}

    Ein Packsattel für die Benutzung an einem Pferd.

    \begin{tabular}{@{}ll@{}}
        \textbf{\trans{weight_label}:} & 4.00 \\
        \textbf{\trans{price_label}:} & 30 \\
        \textbf{\trans{rarity_label}:} & Gewöhnlich \\
        \textbf{\trans{concealment_label}:} & 0 \\
    \end{tabular}

    \wbif{phasesix}{
        \vspace{3mm}
\faIcon{cube}\hspace{0.5em}    }{}


    \subsection*{Eimer}

    Ein 10l Eimer.

    \begin{tabular}{@{}ll@{}}
        \textbf{\trans{weight_label}:} & 0.70 \\
        \textbf{\trans{price_label}:} & 5 \\
        \textbf{\trans{rarity_label}:} & Gewöhnlich \\
        \textbf{\trans{concealment_label}:} & 0 \\
    \end{tabular}

    \wbif{phasesix}{
        \vspace{3mm}
\faIcon{cube}\hspace{0.5em}    }{}


    \subsection*{Glasflasche}

    Eine Glasflasche, die mit allem möglichen gefüllt werden kann.

    \begin{tabular}{@{}ll@{}}
        \textbf{\trans{weight_label}:} & 0.20 \\
        \textbf{\trans{price_label}:} & 5 \\
        \textbf{\trans{rarity_label}:} & Gewöhnlich \\
        \textbf{\trans{concealment_label}:} & 0 \\
    \end{tabular}

    \wbif{phasesix}{
        \vspace{3mm}
\faIcon{cube}\hspace{0.5em}    }{}


    \subsection*{Wasserfaß}

    Dieses Faß kann mit 20l Flüssigkeit gefüllt werden.

    \begin{tabular}{@{}ll@{}}
        \textbf{\trans{weight_label}:} & 5.00 \\
        \textbf{\trans{price_label}:} & 10 \\
        \textbf{\trans{rarity_label}:} & Gewöhnlich \\
        \textbf{\trans{concealment_label}:} & 1 \\
    \end{tabular}

    \wbif{phasesix}{
        \vspace{3mm}
\faIcon{cube}\hspace{0.5em}    }{}


    \subsection*{Jadeschatulle}


    \begin{tabular}{@{}ll@{}}
        \textbf{\trans{weight_label}:} & 0.50 \\
        \textbf{\trans{price_label}:} & 50 \\
        \textbf{\trans{rarity_label}:} & Außergewöhnlich \\
        \textbf{\trans{concealment_label}:} & 0 \\
    \end{tabular}

    \wbif{phasesix}{
        \vspace{3mm}
\faIcon{cube}\hspace{0.5em}    }{}


    \section{Werkzeuge}



    \subsection*{Zirkel}

    Ein Zirkel kann für die Navigation oder geometrische Aufgaben benutzt werden.

    \begin{tabular}{@{}ll@{}}
        \textbf{\trans{weight_label}:} & 0.20 \\
        \textbf{\trans{price_label}:} & 30 \\
        \textbf{\trans{rarity_label}:} & Gewöhnlich \\
        \textbf{\trans{concealment_label}:} & 0 \\
    \end{tabular}

    \wbif{phasesix}{
        \vspace{3mm}
\faIcon{cube}\hspace{0.5em}    }{}


    \subsection*{Ledergebundenes Notizbuch (30 Seiten, A6)}

    Ein kleines Notizbuch, mit einem Schutzeinband aus weichem Leder.

    \begin{tabular}{@{}ll@{}}
        \textbf{\trans{weight_label}:} & 0.20 \\
        \textbf{\trans{price_label}:} & 0 \\
        \textbf{\trans{rarity_label}:} & Gewöhnlich \\
        \textbf{\trans{concealment_label}:} & 0 \\
    \end{tabular}

    \wbif{phasesix}{
        \vspace{3mm}
\faIcon{flag}\hspace{0.5em}\faIcon{plane}\hspace{0.5em}\faIcon{microchip}\hspace{0.5em}\faIcon{fire}\hspace{0.5em}    }{}


    \subsection*{Psychic Paper}

    Auf diesem seltsamen Artefakt einer weit entfernten, lange ausgestorbenen Zivilisation sieht der Betrachter immer das, was der Träger dieses Artefakts wünscht. Es kann jeden erdenklichen
Ausweis oder jedes andere Dokument abbilden und spontan sein Aussehen ändern. Es hat keine Ladung und kann immer wieder verwendet werden.

    \begin{tabular}{@{}ll@{}}
        \textbf{\trans{weight_label}:} & 0.10 \\
        \textbf{\trans{price_label}:} & 0 \\
        \textbf{\trans{rarity_label}:} & Selten \\
        \textbf{\trans{concealment_label}:} & 0 \\
    \end{tabular}

    \wbif{phasesix}{
        \vspace{3mm}
\faIcon{fire}\hspace{0.5em}    }{}


    \subsection*{Mörser und Stößel}


    \begin{tabular}{@{}ll@{}}
        \textbf{\trans{weight_label}:} & 0.50 \\
        \textbf{\trans{price_label}:} & 5 \\
        \textbf{\trans{rarity_label}:} & Gewöhnlich \\
        \textbf{\trans{concealment_label}:} & 0 \\
    \end{tabular}

    \wbif{phasesix}{
        \vspace{3mm}
\faIcon{cube}\hspace{0.5em}    }{}


    \subsection*{Externes Personenschild}

    Dieses Gerät ist ein elektronisches Amulett, welches an einer Kette befestigt ist. Um den 
Hals getragen kann es jederzeit aktiviert werden, wobei die Ladung des Personenschildes 
nur für eine einzige Verwendung reicht. 

Das Personenschild verhindert den kompletten Schaden von drei Angriffen. Verursacht ein Angriff
keinen Schaden so verbraucht er keine Ladung des Schildes. Der Schaden wird nach Deckungs- und 
Schutzwürfen berücksichtigt. 

Das Personenschild kann nur im HQ wieder aufgeladen werden.

    \begin{tabular}{@{}ll@{}}
        \textbf{\trans{weight_label}:} & 0.20 \\
        \textbf{\trans{price_label}:} & 0 \\
        \textbf{\trans{rarity_label}:} & Außergewöhnlich \\
        \textbf{\trans{concealment_label}:} & 0 \\
    \end{tabular}

    \wbif{phasesix}{
        \vspace{3mm}
\faIcon{fire}\hspace{0.5em}    }{}


    \subsection*{Kabelbinder}


    \begin{tabular}{@{}ll@{}}
        \textbf{\trans{weight_label}:} & 0.10 \\
        \textbf{\trans{price_label}:} & 1 \\
        \textbf{\trans{rarity_label}:} & Gewöhnlich \\
        \textbf{\trans{concealment_label}:} & 0 \\
    \end{tabular}

    \wbif{phasesix}{
        \vspace{3mm}
\faIcon{microchip}\hspace{0.5em}    }{}


    \subsection*{Senkblei}

    Ein Senkblei, um etwa die Tiefe von etwas abzuschätzen.

    \begin{tabular}{@{}ll@{}}
        \textbf{\trans{weight_label}:} & 0.30 \\
        \textbf{\trans{price_label}:} & 10 \\
        \textbf{\trans{rarity_label}:} & Gewöhnlich \\
        \textbf{\trans{concealment_label}:} & 0 \\
    \end{tabular}

    \wbif{phasesix}{
        \vspace{3mm}
\faIcon{cube}\hspace{0.5em}    }{}


    \subsection*{Temporalkommunikator}

    Dieses Gerät ermöglicht die Kommunikation aus anderen Zeiten und gar anderen Welten mit dem 
N.E.X.U.S. HQ. Hierdurch sind die Agenten in der Lage einen Rückholauftrag zu senden oder 
einfach nur mit den Mitarbeitern des HQ zu sprechen.

    \begin{tabular}{@{}ll@{}}
        \textbf{\trans{weight_label}:} & 0.50 \\
        \textbf{\trans{price_label}:} & 0 \\
        \textbf{\trans{rarity_label}:} & Selten \\
        \textbf{\trans{concealment_label}:} & 1 \\
    \end{tabular}

    \wbif{phasesix}{
        \vspace{3mm}
\faIcon{fire}\hspace{0.5em}    }{}


    \subsection*{Abakus}

    Der Abakus ist eine einfache Rechenmaschine. Kommt er zur Anwendung sind alle Mechanik Würfe leichte Aürfe.

    \begin{tabular}{@{}ll@{}}
        \textbf{\trans{weight_label}:} & 0.70 \\
        \textbf{\trans{price_label}:} & 80 \\
        \textbf{\trans{rarity_label}:} & Gewöhnlich \\
        \textbf{\trans{concealment_label}:} & 0 \\
    \end{tabular}

    \wbif{phasesix}{
        \vspace{3mm}
\faIcon{cube}\hspace{0.5em}    }{}


    \subsection*{Kleiner Kessel}

    Ein kleiner Kessel aus Eisen

    \begin{tabular}{@{}ll@{}}
        \textbf{\trans{weight_label}:} & 1.00 \\
        \textbf{\trans{price_label}:} & 5 \\
        \textbf{\trans{rarity_label}:} & Gewöhnlich \\
        \textbf{\trans{concealment_label}:} & 0 \\
    \end{tabular}

    \wbif{phasesix}{
        \vspace{3mm}
\faIcon{cube}\hspace{0.5em}    }{}


    \subsection*{Papier}

    Ein Blatt Papier. Man kann unter anderem darauf schreiben.

    \begin{tabular}{@{}ll@{}}
        \textbf{\trans{weight_label}:} & 0.10 \\
        \textbf{\trans{price_label}:} & 1 \\
        \textbf{\trans{rarity_label}:} & Gewöhnlich \\
        \textbf{\trans{concealment_label}:} & 0 \\
    \end{tabular}

    \wbif{phasesix}{
        \vspace{3mm}
\faIcon{umbrella}\hspace{0.5em}\faIcon{flag}\hspace{0.5em}\faIcon{plane}\hspace{0.5em}\faIcon{microchip}\hspace{0.5em}    }{}


    \subsection*{Fesselseil}

    Dieses Fesselseil ist dafür geeignet feste Knoten zu knüpfen.

    \begin{tabular}{@{}ll@{}}
        \textbf{\trans{weight_label}:} & 1.00 \\
        \textbf{\trans{price_label}:} & 20 \\
        \textbf{\trans{rarity_label}:} & Gewöhnlich \\
        \textbf{\trans{concealment_label}:} & 0 \\
    \end{tabular}

    \wbif{phasesix}{
        \vspace{3mm}
\faIcon{cube}\hspace{0.5em}    }{}


    \subsection*{Kleiner Webrahmen}

    Ein kleiner Webrahmen, um auf der Reise gewebte Stoffe herstellen zu können.

    \begin{tabular}{@{}ll@{}}
        \textbf{\trans{weight_label}:} & 2.00 \\
        \textbf{\trans{price_label}:} & 20 \\
        \textbf{\trans{rarity_label}:} & Gewöhnlich \\
        \textbf{\trans{concealment_label}:} & 1 \\
    \end{tabular}

    \wbif{phasesix}{
        \vspace{3mm}
\faIcon{cube}\hspace{0.5em}    }{}


    \subsection*{Brechstange}

    Gordon Freeman knows how to use it

    \begin{tabular}{@{}ll@{}}
        \textbf{\trans{weight_label}:} & 1.00 \\
        \textbf{\trans{price_label}:} & 29 \\
        \textbf{\trans{rarity_label}:} & Gewöhnlich \\
        \textbf{\trans{concealment_label}:} & 1 \\
    \end{tabular}

    \wbif{phasesix}{
        \vspace{3mm}
\faIcon{cube}\hspace{0.5em}    }{}


    \subsection*{Schaufel}


    \begin{tabular}{@{}ll@{}}
        \textbf{\trans{weight_label}:} & 1.00 \\
        \textbf{\trans{price_label}:} & 30 \\
        \textbf{\trans{rarity_label}:} & Gewöhnlich \\
        \textbf{\trans{concealment_label}:} & 3 \\
    \end{tabular}

    \wbif{phasesix}{
        \vspace{3mm}
\faIcon{cube}\hspace{0.5em}    }{}


    \subsection*{Improvisierter Diedrich}


    \begin{tabular}{@{}ll@{}}
        \textbf{\trans{weight_label}:} & 0.01 \\
        \textbf{\trans{price_label}:} & 0 \\
        \textbf{\trans{rarity_label}:} & Gewöhnlich \\
        \textbf{\trans{concealment_label}:} & 0 \\
    \end{tabular}

    \wbif{phasesix}{
        \vspace{3mm}
\faIcon{cube}\hspace{0.5em}    }{}


    \subsection*{Blitzableiter}

    Persönlicher Blitzschutz - In Reihe geschaltet auch als Gruppenschutz in Hochenergie Wäldern

    \begin{tabular}{@{}ll@{}}
        \textbf{\trans{weight_label}:} & 5.00 \\
        \textbf{\trans{price_label}:} & 1000 \\
        \textbf{\trans{rarity_label}:} & Gewöhnlich \\
        \textbf{\trans{concealment_label}:} & 0 \\
    \end{tabular}

    \wbif{phasesix}{
        \vspace{3mm}
\faIcon{space-shuttle}\hspace{0.5em}\faIcon{fire}\hspace{0.5em}    }{}


    \subsection*{Nägel}

    Eine handvoll simpler Nägel

    \begin{tabular}{@{}ll@{}}
        \textbf{\trans{weight_label}:} & 0.05 \\
        \textbf{\trans{price_label}:} & 0 \\
        \textbf{\trans{rarity_label}:} & Gewöhnlich \\
        \textbf{\trans{concealment_label}:} & 0 \\
    \end{tabular}

    \wbif{phasesix}{
        \vspace{3mm}
\faIcon{cube}\hspace{0.5em}    }{}


    \subsection*{Kohlestifte}

    Kohlestifte können verwendet werden um auf Pergament oder Papier zu schreiben.

    \begin{tabular}{@{}ll@{}}
        \textbf{\trans{weight_label}:} & 0.30 \\
        \textbf{\trans{price_label}:} & 5 \\
        \textbf{\trans{rarity_label}:} & Gewöhnlich \\
        \textbf{\trans{concealment_label}:} & 0 \\
    \end{tabular}

    \wbif{phasesix}{
        \vspace{3mm}
\faIcon{cube}\hspace{0.5em}    }{}


    \subsection*{Obsidian Ritualdolch}


    \begin{tabular}{@{}ll@{}}
        \textbf{\trans{weight_label}:} & 1.00 \\
        \textbf{\trans{price_label}:} & 100 \\
        \textbf{\trans{rarity_label}:} & Außergewöhnlich \\
        \textbf{\trans{concealment_label}:} & 0 \\
    \end{tabular}

    \wbif{phasesix}{
        \vspace{3mm}
\faIcon{cube}\hspace{0.5em}    }{}


    \subsection*{Pfeife}

    Eine Pfeife zum Rauchen von Tabak und ähnlichem.

    \begin{tabular}{@{}ll@{}}
        \textbf{\trans{weight_label}:} & 0.10 \\
        \textbf{\trans{price_label}:} & 100 \\
        \textbf{\trans{rarity_label}:} & Gewöhnlich \\
        \textbf{\trans{concealment_label}:} & 0 \\
    \end{tabular}

    \wbif{phasesix}{
        \vspace{3mm}
\faIcon{cube}\hspace{0.5em}    }{}


    \subsection*{Kleine Pfanne}


    \begin{tabular}{@{}ll@{}}
        \textbf{\trans{weight_label}:} & 1.00 \\
        \textbf{\trans{price_label}:} & 5 \\
        \textbf{\trans{rarity_label}:} & Gewöhnlich \\
        \textbf{\trans{concealment_label}:} & 0 \\
    \end{tabular}

    \wbif{phasesix}{
        \vspace{3mm}
\faIcon{cube}\hspace{0.5em}    }{}


    \subsection*{Flaschenzug}

    Ein einfacher Flaschenzug. Ein Seil wird für den Betrieb benötigt. Der Flaschenzug kann 100kg heben.

    \begin{tabular}{@{}ll@{}}
        \textbf{\trans{weight_label}:} & 2.00 \\
        \textbf{\trans{price_label}:} & 40 \\
        \textbf{\trans{rarity_label}:} & Gewöhnlich \\
        \textbf{\trans{concealment_label}:} & 0 \\
    \end{tabular}

    \wbif{phasesix}{
        \vspace{3mm}
\faIcon{cube}\hspace{0.5em}    }{}


    \subsection*{Pinsel}

    Nutze diesen Pinsel um auf einer Leinwand zu malen.

    \begin{tabular}{@{}ll@{}}
        \textbf{\trans{weight_label}:} & 0.10 \\
        \textbf{\trans{price_label}:} & 5 \\
        \textbf{\trans{rarity_label}:} & Gewöhnlich \\
        \textbf{\trans{concealment_label}:} & 0 \\
    \end{tabular}

    \wbif{phasesix}{
        \vspace{3mm}
\faIcon{cube}\hspace{0.5em}    }{}


    \subsection*{Omnisensor}

    Der Omnisensor ist ein Handheld-Gerät mit einem 7 Zoll Display. Das Gerät ist in der Lage die nähere
Umgebung nach dem Vorkommen eines Materials, eines Lebewesens oder einer einprogrammierten
Signatur zu scannen. Zudem zeigt es Daten zur Zusammensetzung der Atmosphäre, der näheren
Umgebung sowie allen Vorkommens von Elektrizität oder Magie in der Nähe an.

    \begin{tabular}{@{}ll@{}}
        \textbf{\trans{weight_label}:} & 0.50 \\
        \textbf{\trans{price_label}:} & 0 \\
        \textbf{\trans{rarity_label}:} & Außergewöhnlich \\
        \textbf{\trans{concealment_label}:} & 2 \\
    \end{tabular}

    \wbif{phasesix}{
        \vspace{3mm}
\faIcon{fire}\hspace{0.5em}    }{}


    \subsection*{Dietriche}

    Wird ein Dietrich mit dem Wissen Schloss knacken verwendet, so wird ein leichter Wurf statt eines normalen Wurfes geworfen.

    \begin{tabular}{@{}ll@{}}
        \textbf{\trans{weight_label}:} & 0.20 \\
        \textbf{\trans{price_label}:} & 30 \\
        \textbf{\trans{rarity_label}:} & Gewöhnlich \\
        \textbf{\trans{concealment_label}:} & 0 \\
    \end{tabular}

    \wbif{phasesix}{
        \vspace{3mm}
\faIcon{cube}\hspace{0.5em}    }{}


    \subsection*{Hammer}


    \begin{tabular}{@{}ll@{}}
        \textbf{\trans{weight_label}:} & 2.00 \\
        \textbf{\trans{price_label}:} & 30 \\
        \textbf{\trans{rarity_label}:} & Gewöhnlich \\
        \textbf{\trans{concealment_label}:} & 1 \\
    \end{tabular}

    \wbif{phasesix}{
        \vspace{3mm}
\faIcon{cube}\hspace{0.5em}    }{}


    \subsection*{Reisigbesen}

    Ein Besen. Man kann mit ihm fegen.

    \begin{tabular}{@{}ll@{}}
        \textbf{\trans{weight_label}:} & 2.00 \\
        \textbf{\trans{price_label}:} & 10 \\
        \textbf{\trans{rarity_label}:} & Gewöhnlich \\
        \textbf{\trans{concealment_label}:} & 1 \\
    \end{tabular}

    \wbif{phasesix}{
        \vspace{3mm}
\faIcon{cube}\hspace{0.5em}    }{}


    \subsection*{Fusionszelle}

    Betreibt futuristische Gegenstände

    \begin{tabular}{@{}ll@{}}
        \textbf{\trans{weight_label}:} & 0.50 \\
        \textbf{\trans{price_label}:} & 500 \\
        \textbf{\trans{rarity_label}:} & Selten \\
        \textbf{\trans{concealment_label}:} & 2 \\
        \textbf{\trans{charges_label}:} & 1\\
    \end{tabular}

    \wbif{phasesix}{
        \vspace{3mm}
\faIcon{space-shuttle}\hspace{0.5em}\faIcon{fire}\hspace{0.5em}    }{}


    \subsection*{Schiefertafel}

    Auf dieser Schiefertafel kann man schreiben, und man kann das geschriebene jederzeit wegwischen.

    \begin{tabular}{@{}ll@{}}
        \textbf{\trans{weight_label}:} & 0.50 \\
        \textbf{\trans{price_label}:} & 10 \\
        \textbf{\trans{rarity_label}:} & Gewöhnlich \\
        \textbf{\trans{concealment_label}:} & 0 \\
    \end{tabular}

    \wbif{phasesix}{
        \vspace{3mm}
\faIcon{cube}\hspace{0.5em}    }{}


    \subsection*{Tintenfass}

    Ein sicher verschlossenes Tintenfass, welches Tinte für eine Feder oder einen Gänsekiel enthält.

    \begin{tabular}{@{}ll@{}}
        \textbf{\trans{weight_label}:} & 0.60 \\
        \textbf{\trans{price_label}:} & 10 \\
        \textbf{\trans{rarity_label}:} & Gewöhnlich \\
        \textbf{\trans{concealment_label}:} & 0 \\
        \textbf{\trans{charges_label}:} & 25\\
    \end{tabular}

    \wbif{phasesix}{
        \vspace{3mm}
\faIcon{cube}\hspace{0.5em}    }{}


    \subsection*{Feuerzeug}


    \begin{tabular}{@{}ll@{}}
        \textbf{\trans{weight_label}:} & 0.10 \\
        \textbf{\trans{price_label}:} & 1 \\
        \textbf{\trans{rarity_label}:} & Gewöhnlich \\
        \textbf{\trans{concealment_label}:} & 0 \\
    \end{tabular}

    \wbif{phasesix}{
        \vspace{3mm}
\faIcon{flag}\hspace{0.5em}\faIcon{plane}\hspace{0.5em}\faIcon{microchip}\hspace{0.5em}    }{}


    \section{Lichter}



    \subsection*{Taschenlampe}


    \begin{tabular}{@{}ll@{}}
        \textbf{\trans{weight_label}:} & 0.80 \\
        \textbf{\trans{price_label}:} & 40 \\
        \textbf{\trans{rarity_label}:} & Gewöhnlich \\
        \textbf{\trans{concealment_label}:} & 1 \\
    \end{tabular}

    \wbif{phasesix}{
        \vspace{3mm}
\faIcon{flag}\hspace{0.5em}\faIcon{plane}\hspace{0.5em}\faIcon{microchip}\hspace{0.5em}    }{}


    \subsection*{Fackel}


    \begin{tabular}{@{}ll@{}}
        \textbf{\trans{weight_label}:} & 0.20 \\
        \textbf{\trans{price_label}:} & 2 \\
        \textbf{\trans{rarity_label}:} & Gewöhnlich \\
        \textbf{\trans{concealment_label}:} & 0 \\
    \end{tabular}

    \wbif{phasesix}{
        \vspace{3mm}
\faIcon{cube}\hspace{0.5em}    }{}


    \subsection*{Laterne}


    \begin{tabular}{@{}ll@{}}
        \textbf{\trans{weight_label}:} & 1.00 \\
        \textbf{\trans{price_label}:} & 40 \\
        \textbf{\trans{rarity_label}:} & Gewöhnlich \\
        \textbf{\trans{concealment_label}:} & 1 \\
    \end{tabular}

    \wbif{phasesix}{
        \vspace{3mm}
\faIcon{cube}\hspace{0.5em}    }{}


    \subsection*{Magnesium Fackel}

    Eine brennende, helle Fackel. Taucht die Umgebung in helles, rötliches Licht.

    \begin{tabular}{@{}ll@{}}
        \textbf{\trans{weight_label}:} & 0.20 \\
        \textbf{\trans{price_label}:} & 15 \\
        \textbf{\trans{rarity_label}:} & Gewöhnlich \\
        \textbf{\trans{concealment_label}:} & 0 \\
    \end{tabular}

    \wbif{phasesix}{
        \vspace{3mm}
\faIcon{flag}\hspace{0.5em}\faIcon{plane}\hspace{0.5em}\faIcon{microchip}\hspace{0.5em}\faIcon{space-shuttle}\hspace{0.5em}    }{}


    \subsection*{Kerze}

    Eine Kerze. Brennt etwa 8 Stunden.

    \begin{tabular}{@{}ll@{}}
        \textbf{\trans{weight_label}:} & 0.20 \\
        \textbf{\trans{price_label}:} & 5 \\
        \textbf{\trans{rarity_label}:} & Gewöhnlich \\
        \textbf{\trans{concealment_label}:} & 0 \\
    \end{tabular}

    \wbif{phasesix}{
        \vspace{3mm}
\faIcon{cube}\hspace{0.5em}    }{}


    \subsection*{Pechfackel}

    Die Pechfackel brennt für etwa 8 Stunden und erzeugt ein angenehmes, großflächiges Licht.

    \begin{tabular}{@{}ll@{}}
        \textbf{\trans{weight_label}:} & 0.50 \\
        \textbf{\trans{price_label}:} & 10 \\
        \textbf{\trans{rarity_label}:} & Gewöhnlich \\
        \textbf{\trans{concealment_label}:} & 0 \\
    \end{tabular}

    \wbif{phasesix}{
        \vspace{3mm}
\faIcon{cube}\hspace{0.5em}    }{}


    \subsection*{Öllampe}

    Die Öllampe verbreitet ein angenehmes, großflächiges Licht, und ist nicht so windanfällig wie eine Fackel.

    \begin{tabular}{@{}ll@{}}
        \textbf{\trans{weight_label}:} & 1.00 \\
        \textbf{\trans{price_label}:} & 30 \\
        \textbf{\trans{rarity_label}:} & Gewöhnlich \\
        \textbf{\trans{concealment_label}:} & 0 \\
    \end{tabular}

    \wbif{phasesix}{
        \vspace{3mm}
\faIcon{cube}\hspace{0.5em}    }{}


    \subsection*{Sturmlaterne}

    Die Sturmlaterne ist besonders resistent gegen Wind und Wetter. Sie verbreitet ein angenehmes Licht.

    \begin{tabular}{@{}ll@{}}
        \textbf{\trans{weight_label}:} & 1.00 \\
        \textbf{\trans{price_label}:} & 60 \\
        \textbf{\trans{rarity_label}:} & Gewöhnlich \\
        \textbf{\trans{concealment_label}:} & 0 \\
    \end{tabular}

    \wbif{phasesix}{
        \vspace{3mm}
\faIcon{cube}\hspace{0.5em}    }{}


    \section{Überwachung}



    \subsection*{Kamera}


    \begin{tabular}{@{}ll@{}}
        \textbf{\trans{weight_label}:} & 0.80 \\
        \textbf{\trans{price_label}:} & 100 \\
        \textbf{\trans{rarity_label}:} & Gewöhnlich \\
        \textbf{\trans{concealment_label}:} & 1 \\
    \end{tabular}

    \wbif{phasesix}{
        \vspace{3mm}
\faIcon{flag}\hspace{0.5em}\faIcon{plane}\hspace{0.5em}\faIcon{microchip}\hspace{0.5em}    }{}


    \subsection*{Digitalkamera}


    \begin{tabular}{@{}ll@{}}
        \textbf{\trans{weight_label}:} & 0.80 \\
        \textbf{\trans{price_label}:} & 400 \\
        \textbf{\trans{rarity_label}:} & Gewöhnlich \\
        \textbf{\trans{concealment_label}:} & 2 \\
    \end{tabular}

    \wbif{phasesix}{
        \vspace{3mm}
\faIcon{microchip}\hspace{0.5em}\faIcon{space-shuttle}\hspace{0.5em}    }{}


    \subsection*{Fernrohr}

    Alle \textit{Wahrnehmungswürfe}, die mit Hilfe des Fernrohrs gemacht werden sind einfache Proben.

    \begin{tabular}{@{}ll@{}}
        \textbf{\trans{weight_label}:} & 0.50 \\
        \textbf{\trans{price_label}:} & 80 \\
        \textbf{\trans{rarity_label}:} & Gewöhnlich \\
        \textbf{\trans{concealment_label}:} & 0 \\
    \end{tabular}

    \wbif{phasesix}{
        \vspace{3mm}
\faIcon{cube}\hspace{0.5em}    }{}


    \subsection*{Wanze (Abhörgerät)}

    Kann zum Abhören von Personen und Räumen plaziert werden.

    \begin{tabular}{@{}ll@{}}
        \textbf{\trans{weight_label}:} & 0.00 \\
        \textbf{\trans{price_label}:} & 150 \\
        \textbf{\trans{rarity_label}:} & Außergewöhnlich \\
        \textbf{\trans{concealment_label}:} & 1 \\
    \end{tabular}

    \wbif{phasesix}{
        \vspace{3mm}
\faIcon{plane}\hspace{0.5em}\faIcon{microchip}\hspace{0.5em}    }{}


    \subsection*{Handschnellen}


    \begin{tabular}{@{}ll@{}}
        \textbf{\trans{weight_label}:} & 0.30 \\
        \textbf{\trans{price_label}:} & 50 \\
        \textbf{\trans{rarity_label}:} & Gewöhnlich \\
        \textbf{\trans{concealment_label}:} & 1 \\
    \end{tabular}

    \wbif{phasesix}{
        \vspace{3mm}
\faIcon{flag}\hspace{0.5em}\faIcon{plane}\hspace{0.5em}\faIcon{microchip}\hspace{0.5em}\faIcon{space-shuttle}\hspace{0.5em}    }{}


    \section{Kommunikation}



    \subsection*{Handy}


    \begin{tabular}{@{}ll@{}}
        \textbf{\trans{weight_label}:} & 0.08 \\
        \textbf{\trans{price_label}:} & 100 \\
        \textbf{\trans{rarity_label}:} & Gewöhnlich \\
        \textbf{\trans{concealment_label}:} & 0 \\
    \end{tabular}

    \wbif{phasesix}{
        \vspace{3mm}
\faIcon{microchip}\hspace{0.5em}    }{}


    \subsection*{Smartphone}


    \begin{tabular}{@{}ll@{}}
        \textbf{\trans{weight_label}:} & 0.20 \\
        \textbf{\trans{price_label}:} & 500 \\
        \textbf{\trans{rarity_label}:} & Gewöhnlich \\
        \textbf{\trans{concealment_label}:} & 0 \\
    \end{tabular}

    \wbif{phasesix}{
        \vspace{3mm}
\faIcon{microchip}\hspace{0.5em}    }{}


    \subsection*{Signalpfeife}


    \begin{tabular}{@{}ll@{}}
        \textbf{\trans{weight_label}:} & 0.03 \\
        \textbf{\trans{price_label}:} & 17 \\
        \textbf{\trans{rarity_label}:} & Gewöhnlich \\
        \textbf{\trans{concealment_label}:} & 0 \\
    \end{tabular}

    \wbif{phasesix}{
        \vspace{3mm}
\faIcon{umbrella}\hspace{0.5em}\faIcon{flag}\hspace{0.5em}\faIcon{plane}\hspace{0.5em}\faIcon{microchip}\hspace{0.5em}    }{}


    \subsection*{Universalkommunikator}

    Der Universalkommunikator ist in der Lage Sprache in nahezu alle Dialekte des Universums zu 
übersetzen. Ein Agent der den Universalkommunikator trägt kann sich mit jedem intelligenten
Wesen unterhalten.

    \begin{tabular}{@{}ll@{}}
        \textbf{\trans{weight_label}:} & 0.50 \\
        \textbf{\trans{price_label}:} & 0 \\
        \textbf{\trans{rarity_label}:} & Außergewöhnlich \\
        \textbf{\trans{concealment_label}:} & 1 \\
    \end{tabular}

    \wbif{phasesix}{
        \vspace{3mm}
\faIcon{fire}\hspace{0.5em}    }{}


    \subsection*{Intercom-Funk}


    \begin{tabular}{@{}ll@{}}
        \textbf{\trans{weight_label}:} & 0.10 \\
        \textbf{\trans{price_label}:} & 250 \\
        \textbf{\trans{rarity_label}:} & Gewöhnlich \\
        \textbf{\trans{concealment_label}:} & 0 \\
    \end{tabular}

    \wbif{phasesix}{
        \vspace{3mm}
\faIcon{microchip}\hspace{0.5em}    }{}


    \subsection*{USB Stick}

    Ein einfacher USB Stick.

    \begin{tabular}{@{}ll@{}}
        \textbf{\trans{weight_label}:} & 0.02 \\
        \textbf{\trans{price_label}:} & 10 \\
        \textbf{\trans{rarity_label}:} & Gewöhnlich \\
        \textbf{\trans{concealment_label}:} & 0 \\
    \end{tabular}

    \wbif{phasesix}{
        \vspace{3mm}
\faIcon{microchip}\hspace{0.5em}    }{}


    \subsection*{Bankkarte}

    Eine Bankkarte aus Plastik, wahlweise mit EC oder Kreditkartenfunktion.

    \begin{tabular}{@{}ll@{}}
        \textbf{\trans{weight_label}:} & 0.01 \\
        \textbf{\trans{price_label}:} & 10 \\
        \textbf{\trans{rarity_label}:} & Gewöhnlich \\
        \textbf{\trans{concealment_label}:} & 0 \\
    \end{tabular}

    \wbif{phasesix}{
        \vspace{3mm}
\faIcon{plane}\hspace{0.5em}\faIcon{microchip}\hspace{0.5em}\faIcon{space-shuttle}\hspace{0.5em}    }{}


    \section{Trekking Ausrüstung}



    \subsection*{Angelhaken und Schnur}

    Eine simple Anlgerausrüstung.

    \begin{tabular}{@{}ll@{}}
        \textbf{\trans{weight_label}:} & 0.20 \\
        \textbf{\trans{price_label}:} & 10 \\
        \textbf{\trans{rarity_label}:} & Gewöhnlich \\
        \textbf{\trans{concealment_label}:} & 0 \\
    \end{tabular}

    \wbif{phasesix}{
        \vspace{3mm}
\faIcon{cube}\hspace{0.5em}    }{}


    \subsection*{Wasserschlauch}

    Ein 1 Liter Lederschlauch um Wasser zu transportieren.

    \begin{tabular}{@{}ll@{}}
        \textbf{\trans{weight_label}:} & 0.30 \\
        \textbf{\trans{price_label}:} & 20 \\
        \textbf{\trans{rarity_label}:} & Gewöhnlich \\
        \textbf{\trans{concealment_label}:} & 0 \\
    \end{tabular}

    \wbif{phasesix}{
        \vspace{3mm}
\faIcon{cube}\hspace{0.5em}    }{}


    \subsection*{Zelt}

    Ein großes 4-personen Zelt. Es ist ein wenig anstrengend aufzubauen, bietet aber Platz und Schutz für 4-5 Personen.

    \begin{tabular}{@{}ll@{}}
        \textbf{\trans{weight_label}:} & 5.00 \\
        \textbf{\trans{price_label}:} & 70 \\
        \textbf{\trans{rarity_label}:} & Gewöhnlich \\
        \textbf{\trans{concealment_label}:} & 1 \\
    \end{tabular}

    \wbif{phasesix}{
        \vspace{3mm}
\faIcon{cube}\hspace{0.5em}    }{}


    \subsection*{Hängematte}

    Diese Hängematte kann aufgespannt werden um eine bequeme Schlafstätte zu bieten.

    \begin{tabular}{@{}ll@{}}
        \textbf{\trans{weight_label}:} & 2.00 \\
        \textbf{\trans{price_label}:} & 20 \\
        \textbf{\trans{rarity_label}:} & Gewöhnlich \\
        \textbf{\trans{concealment_label}:} & 0 \\
    \end{tabular}

    \wbif{phasesix}{
        \vspace{3mm}
\faIcon{cube}\hspace{0.5em}    }{}


    \subsection*{Seil (3m)}


    \begin{tabular}{@{}ll@{}}
        \textbf{\trans{weight_label}:} & 3.00 \\
        \textbf{\trans{price_label}:} & 30 \\
        \textbf{\trans{rarity_label}:} & Gewöhnlich \\
        \textbf{\trans{concealment_label}:} & 2 \\
    \end{tabular}

    \wbif{phasesix}{
        \vspace{3mm}
\faIcon{cube}\hspace{0.5em}    }{}


    \subsection*{Wurfhaken}

    Ein Wurfhaken, dazu gedacht dorthin geworfen zu werden wo er sich festhaken kann. Idealerweise wird er zusammen mit einem daran festgebundenen Seil verwendet.

    \begin{tabular}{@{}ll@{}}
        \textbf{\trans{weight_label}:} & 2.00 \\
        \textbf{\trans{price_label}:} & 90 \\
        \textbf{\trans{rarity_label}:} & Gewöhnlich \\
        \textbf{\trans{concealment_label}:} & 1 \\
    \end{tabular}

    \wbif{phasesix}{
        \vspace{3mm}
\faIcon{cube}\hspace{0.5em}    }{}


    \subsection*{Rucksack}


    \begin{tabular}{@{}ll@{}}
        \textbf{\trans{weight_label}:} & 1.20 \\
        \textbf{\trans{price_label}:} & 100 \\
        \textbf{\trans{rarity_label}:} & Gewöhnlich \\
        \textbf{\trans{concealment_label}:} & 2 \\
    \end{tabular}

    \wbif{phasesix}{
        \vspace{3mm}
\faIcon{cube}\hspace{0.5em}    }{}


    \subsection*{Feuerstein und Stahl}

    Eine Möglichkeit ein Feuer zu entfachen. Ein wenig anstrengend, aber eine sehr sichere Methode.

    \begin{tabular}{@{}ll@{}}
        \textbf{\trans{weight_label}:} & 0.20 \\
        \textbf{\trans{price_label}:} & 5 \\
        \textbf{\trans{rarity_label}:} & Gewöhnlich \\
        \textbf{\trans{concealment_label}:} & 0 \\
    \end{tabular}

    \wbif{phasesix}{
        \vspace{3mm}
\faIcon{cube}\hspace{0.5em}    }{}


    \subsection*{Decke}


    \begin{tabular}{@{}ll@{}}
        \textbf{\trans{weight_label}:} & 1.00 \\
        \textbf{\trans{price_label}:} & 50 \\
        \textbf{\trans{rarity_label}:} & Gewöhnlich \\
        \textbf{\trans{concealment_label}:} & 1 \\
    \end{tabular}

    \wbif{phasesix}{
        \vspace{3mm}
\faIcon{cube}\hspace{0.5em}    }{}


    \subsection*{Gürteltaschen}

    Bequem zu erreichende Gürteltaschen. Etwa 4 davon können an einem Gürtel angebracht werden.

    \begin{tabular}{@{}ll@{}}
        \textbf{\trans{weight_label}:} & 0.30 \\
        \textbf{\trans{price_label}:} & 30 \\
        \textbf{\trans{rarity_label}:} & Gewöhnlich \\
        \textbf{\trans{concealment_label}:} & 0 \\
    \end{tabular}

    \wbif{phasesix}{
        \vspace{3mm}
\faIcon{cube}\hspace{0.5em}    }{}


    \subsection*{Brennglas}

    Eine Lupe, mit der unter anderem ein Feuer entzündet werden kann.

    \begin{tabular}{@{}ll@{}}
        \textbf{\trans{weight_label}:} & 0.20 \\
        \textbf{\trans{price_label}:} & 50 \\
        \textbf{\trans{rarity_label}:} & Gewöhnlich \\
        \textbf{\trans{concealment_label}:} & 0 \\
    \end{tabular}

    \wbif{phasesix}{
        \vspace{3mm}
\faIcon{cube}\hspace{0.5em}    }{}


    \subsection*{Lampenöl}

    Ein Gefäß voller Lampenöl, um Sturmlaternen oder Öllampen nachzufüllen.

    \begin{tabular}{@{}ll@{}}
        \textbf{\trans{weight_label}:} & 1.00 \\
        \textbf{\trans{price_label}:} & 20 \\
        \textbf{\trans{rarity_label}:} & Gewöhnlich \\
        \textbf{\trans{concealment_label}:} & 0 \\
    \end{tabular}

    \wbif{phasesix}{
        \vspace{3mm}
\faIcon{cube}\hspace{0.5em}    }{}


    \subsection*{Ski}

    Ein paar Ski, die benutzt werden können um sich schnell auf Schnee zu bewegen.

    \begin{tabular}{@{}ll@{}}
        \textbf{\trans{weight_label}:} & 3.00 \\
        \textbf{\trans{price_label}:} & 70 \\
        \textbf{\trans{rarity_label}:} & Gewöhnlich \\
        \textbf{\trans{concealment_label}:} & 0 \\
    \end{tabular}

    \wbif{phasesix}{
        \vspace{3mm}
\faIcon{flag}\hspace{0.5em}\faIcon{plane}\hspace{0.5em}\faIcon{microchip}\hspace{0.5em}\faIcon{space-shuttle}\hspace{0.5em}    }{}


    \subsection*{Strickleiter}

    Wenn die Strickleiter zusammengelegt ist, ist sie einfach zu verstauen. Ausgerollt bietet sie eine spontane Leiter über 8 Meter Höhe.

    \begin{tabular}{@{}ll@{}}
        \textbf{\trans{weight_label}:} & 2.00 \\
        \textbf{\trans{price_label}:} & 40 \\
        \textbf{\trans{rarity_label}:} & Gewöhnlich \\
        \textbf{\trans{concealment_label}:} & 0 \\
    \end{tabular}

    \wbif{phasesix}{
        \vspace{3mm}
\faIcon{cube}\hspace{0.5em}    }{}


    \subsection*{Lasso}

    Dieses Seil ist dafür gemacht ein Lasso zu knüpfen, um Tiere einzufangen.

    \begin{tabular}{@{}ll@{}}
        \textbf{\trans{weight_label}:} & 2.00 \\
        \textbf{\trans{price_label}:} & 20 \\
        \textbf{\trans{rarity_label}:} & Gewöhnlich \\
        \textbf{\trans{concealment_label}:} & 0 \\
    \end{tabular}

    \wbif{phasesix}{
        \vspace{3mm}
\faIcon{cube}\hspace{0.5em}    }{}


    \subsection*{Kompass}

    Zeigt nach Norden

    \begin{tabular}{@{}ll@{}}
        \textbf{\trans{weight_label}:} & 0.10 \\
        \textbf{\trans{price_label}:} & 20 \\
        \textbf{\trans{rarity_label}:} & Gewöhnlich \\
        \textbf{\trans{concealment_label}:} & 0 \\
    \end{tabular}

    \wbif{phasesix}{
        \vspace{3mm}
\faIcon{umbrella}\hspace{0.5em}\faIcon{flag}\hspace{0.5em}\faIcon{plane}\hspace{0.5em}\faIcon{microchip}\hspace{0.5em}    }{}


    \subsection*{Fischnetz}

    Mit diesem Netz kann man gut angeln.

    \begin{tabular}{@{}ll@{}}
        \textbf{\trans{weight_label}:} & 1.00 \\
        \textbf{\trans{price_label}:} & 10 \\
        \textbf{\trans{rarity_label}:} & Gewöhnlich \\
        \textbf{\trans{concealment_label}:} & 0 \\
    \end{tabular}

    \wbif{phasesix}{
        \vspace{3mm}
\faIcon{cube}\hspace{0.5em}    }{}


    \subsection*{Schlafsack}


    \begin{tabular}{@{}ll@{}}
        \textbf{\trans{weight_label}:} & 1.00 \\
        \textbf{\trans{price_label}:} & 50 \\
        \textbf{\trans{rarity_label}:} & Gewöhnlich \\
        \textbf{\trans{concealment_label}:} & 2 \\
    \end{tabular}

    \wbif{phasesix}{
        \vspace{3mm}
\faIcon{cube}\hspace{0.5em}    }{}


    \subsection*{Schneeschuhe}

    Dieses Paar Schneeschuhe kann benutzt werden um auf Schnee bequem und schnell zu laufen.

    \begin{tabular}{@{}ll@{}}
        \textbf{\trans{weight_label}:} & 1.00 \\
        \textbf{\trans{price_label}:} & 20 \\
        \textbf{\trans{rarity_label}:} & Gewöhnlich \\
        \textbf{\trans{concealment_label}:} & 0 \\
    \end{tabular}

    \wbif{phasesix}{
        \vspace{3mm}
\faIcon{cube}\hspace{0.5em}    }{}


    \subsection*{Trockenfleisch}

    Trockenfleisch ist durch Lufttrocknung haltbar gemachtes Fleisch, das aus rohem oder erhitztem Fleisch oder Fleischerzeugnissen hergestellt werden kann.

    \begin{tabular}{@{}ll@{}}
        \textbf{\trans{weight_label}:} & 0.10 \\
        \textbf{\trans{price_label}:} & 5 \\
        \textbf{\trans{rarity_label}:} & Gewöhnlich \\
        \textbf{\trans{concealment_label}:} & 0 \\
    \end{tabular}

    \wbif{phasesix}{
        \vspace{3mm}
\faIcon{cube}\hspace{0.5em}    }{}


    \subsection*{Kletterhaken}

    Ein Kletterhaken kann angebracht werden um Seile darin zu befestigen. Um ihn in den Fels zu schlagen kann ein Hammer verwendet werden.

    \begin{tabular}{@{}ll@{}}
        \textbf{\trans{weight_label}:} & 1.00 \\
        \textbf{\trans{price_label}:} & 5 \\
        \textbf{\trans{rarity_label}:} & Gewöhnlich \\
        \textbf{\trans{concealment_label}:} & 0 \\
    \end{tabular}

    \wbif{phasesix}{
        \vspace{3mm}
\faIcon{cube}\hspace{0.5em}    }{}


    \subsection*{Wurfzelt}


    \begin{tabular}{@{}ll@{}}
        \textbf{\trans{weight_label}:} & 3.00 \\
        \textbf{\trans{price_label}:} & 100 \\
        \textbf{\trans{rarity_label}:} & Gewöhnlich \\
        \textbf{\trans{concealment_label}:} & 5 \\
    \end{tabular}

    \wbif{phasesix}{
        \vspace{3mm}
\faIcon{plane}\hspace{0.5em}\faIcon{microchip}\hspace{0.5em}    }{}


    \section{Nahrung / Proviant}

    Nahrungsgegenstände zur versorgung der Hungrigen Helden und Heldinen


    \subsection*{Italienischer Kräuterlikör 30\% 0,7L}


    \begin{tabular}{@{}ll@{}}
        \textbf{\trans{weight_label}:} & 1.00 \\
        \textbf{\trans{price_label}:} & 16 \\
        \textbf{\trans{rarity_label}:} & Außergewöhnlich \\
        \textbf{\trans{concealment_label}:} & 0 \\
    \end{tabular}

    \wbif{phasesix}{
        \vspace{3mm}
\faIcon{plane}\hspace{0.5em}\faIcon{microchip}\hspace{0.5em}\faIcon{space-shuttle}\hspace{0.5em}    }{}


    \subsection*{Tabak}

    Bestes Langbodenblatt, grobschnitt, vollmundig.

    \begin{tabular}{@{}ll@{}}
        \textbf{\trans{weight_label}:} & 0.05 \\
        \textbf{\trans{price_label}:} & 15 \\
        \textbf{\trans{rarity_label}:} & Gewöhnlich \\
        \textbf{\trans{concealment_label}:} & 0 \\
        \textbf{\trans{charges_label}:} & 20\\
    \end{tabular}

    \wbif{phasesix}{
        \vspace{3mm}
\faIcon{cube}\hspace{0.5em}    }{}


    \subsection*{Edler Wein}

    Eine Flasche guten Weins.

    \begin{tabular}{@{}ll@{}}
        \textbf{\trans{weight_label}:} & 1.00 \\
        \textbf{\trans{price_label}:} & 80 \\
        \textbf{\trans{rarity_label}:} & Gewöhnlich \\
        \textbf{\trans{concealment_label}:} & 0 \\
        \textbf{\trans{charges_label}:} & 3\\
    \end{tabular}

    \wbif{phasesix}{
        \vspace{3mm}
\faIcon{cube}\hspace{0.5em}    }{}


    \subsection*{Bier}

    Kalt, kühl, lecker! Son frisches Bier, Junge, lecker. Kalt muss es sein, Junge!

    \begin{tabular}{@{}ll@{}}
        \textbf{\trans{weight_label}:} & 1.00 \\
        \textbf{\trans{price_label}:} & 1 \\
        \textbf{\trans{rarity_label}:} & Gewöhnlich \\
        \textbf{\trans{concealment_label}:} & 0 \\
    \end{tabular}

    \wbif{phasesix}{
        \vspace{3mm}
\faIcon{cube}\hspace{0.5em}    }{}


    \subsection*{Dörrfleisch}

    Trockenfleisch, nahrhaft und lange haltbar

    \begin{tabular}{@{}ll@{}}
        \textbf{\trans{weight_label}:} & 0.50 \\
        \textbf{\trans{price_label}:} & 1 \\
        \textbf{\trans{rarity_label}:} & Gewöhnlich \\
        \textbf{\trans{concealment_label}:} & 0 \\
        \textbf{\trans{charges_label}:} & 3\\
    \end{tabular}

    \wbif{phasesix}{
        \vspace{3mm}
\faIcon{cube}\hspace{0.5em}    }{}


    \subsection*{Eintopf}

    Eintopf aus verschiedenen Zutaten, alles, was der Koch gefunden hat. Er ist vielleicht etwas schwer zu transportieren, aber der Eintopf enthält sicher viele nahrhafte Zutaten.

    \begin{tabular}{@{}ll@{}}
        \textbf{\trans{weight_label}:} & 0.30 \\
        \textbf{\trans{price_label}:} & 5 \\
        \textbf{\trans{rarity_label}:} & Gewöhnlich \\
        \textbf{\trans{concealment_label}:} & 0 \\
    \end{tabular}

    \wbif{phasesix}{
        \vspace{3mm}
\faIcon{cube}\hspace{0.5em}    }{}


    \section{Fahrzeuge}



    \subsection*{Kanu}

    Das Kanu kann genutzt werden um Wasser zu überqueren. Es ist jedoch nicht hochseetauglich.

    \begin{tabular}{@{}ll@{}}
        \textbf{\trans{weight_label}:} & 20.00 \\
        \textbf{\trans{price_label}:} & 60 \\
        \textbf{\trans{rarity_label}:} & Gewöhnlich \\
        \textbf{\trans{concealment_label}:} & 8 \\
    \end{tabular}

    \wbif{phasesix}{
        \vspace{3mm}
\faIcon{cube}\hspace{0.5em}    }{}


    \subsection*{Kleines Ruderboot}

    Ein Ruderboot samt Rudern.

    \begin{tabular}{@{}ll@{}}
        \textbf{\trans{weight_label}:} & 100.00 \\
        \textbf{\trans{price_label}:} & 120 \\
        \textbf{\trans{rarity_label}:} & Gewöhnlich \\
        \textbf{\trans{concealment_label}:} & 8 \\
    \end{tabular}

    \wbif{phasesix}{
        \vspace{3mm}
\faIcon{cube}\hspace{0.5em}    }{}


    \section{Tierbedarf}



    \subsection*{Pferdefutter}

    Hochwertiges Pferdefutter, eine Portion reicht etwa für eine Woche

    \begin{tabular}{@{}ll@{}}
        \textbf{\trans{weight_label}:} & 1.00 \\
        \textbf{\trans{price_label}:} & 2 \\
        \textbf{\trans{rarity_label}:} & Gewöhnlich \\
        \textbf{\trans{concealment_label}:} & 0 \\
    \end{tabular}

    \wbif{phasesix}{
        \vspace{3mm}
\faIcon{cube}\hspace{0.5em}    }{}


    \subsection*{Tierfutter}

    Hochwertiges Tierfutter. Eine Portion reicht etwa eine Woche.

    \begin{tabular}{@{}ll@{}}
        \textbf{\trans{weight_label}:} & 1.00 \\
        \textbf{\trans{price_label}:} & 1 \\
        \textbf{\trans{rarity_label}:} & Gewöhnlich \\
        \textbf{\trans{concealment_label}:} & 0 \\
    \end{tabular}

    \wbif{phasesix}{
        \vspace{3mm}
\faIcon{cube}\hspace{0.5em}    }{}


    \subsection*{Zaumzeug}


    \begin{tabular}{@{}ll@{}}
        \textbf{\trans{weight_label}:} & 1.00 \\
        \textbf{\trans{price_label}:} & 70 \\
        \textbf{\trans{rarity_label}:} & Gewöhnlich \\
        \textbf{\trans{concealment_label}:} & 0 \\
    \end{tabular}

    \wbif{phasesix}{
        \vspace{3mm}
\faIcon{cube}\hspace{0.5em}    }{}


    \subsection*{Kummet}

    Ein gepolsterter Ring, welcher verwendet wird um Ochsen anzuschirren.

    \begin{tabular}{@{}ll@{}}
        \textbf{\trans{weight_label}:} & 1.00 \\
        \textbf{\trans{price_label}:} & 20 \\
        \textbf{\trans{rarity_label}:} & Gewöhnlich \\
        \textbf{\trans{concealment_label}:} & 0 \\
    \end{tabular}

    \wbif{phasesix}{
        \vspace{3mm}
\faIcon{cube}\hspace{0.5em}    }{}


    \subsection*{Pferdedecke}


    \begin{tabular}{@{}ll@{}}
        \textbf{\trans{weight_label}:} & 2.00 \\
        \textbf{\trans{price_label}:} & 40 \\
        \textbf{\trans{rarity_label}:} & Gewöhnlich \\
        \textbf{\trans{concealment_label}:} & 0 \\
    \end{tabular}

    \wbif{phasesix}{
        \vspace{3mm}
\faIcon{cube}\hspace{0.5em}    }{}


    \subsection*{Sattel}


    \begin{tabular}{@{}ll@{}}
        \textbf{\trans{weight_label}:} & 4.00 \\
        \textbf{\trans{price_label}:} & 80 \\
        \textbf{\trans{rarity_label}:} & Gewöhnlich \\
        \textbf{\trans{concealment_label}:} & 4 \\
    \end{tabular}

    \wbif{phasesix}{
        \vspace{3mm}
\faIcon{cube}\hspace{0.5em}    }{}


    \subsection*{Packsattel}

    Ein Sattel mit Taschen.

    \begin{tabular}{@{}ll@{}}
        \textbf{\trans{weight_label}:} & 5.00 \\
        \textbf{\trans{price_label}:} & 50 \\
        \textbf{\trans{rarity_label}:} & Gewöhnlich \\
        \textbf{\trans{concealment_label}:} & 4 \\
    \end{tabular}

    \wbif{phasesix}{
        \vspace{3mm}
\faIcon{cube}\hspace{0.5em}    }{}


    \subsection*{Striegelzeug}


    \begin{tabular}{@{}ll@{}}
        \textbf{\trans{weight_label}:} & 1.00 \\
        \textbf{\trans{price_label}:} & 30 \\
        \textbf{\trans{rarity_label}:} & Gewöhnlich \\
        \textbf{\trans{concealment_label}:} & 0 \\
    \end{tabular}

    \wbif{phasesix}{
        \vspace{3mm}
\faIcon{cube}\hspace{0.5em}    }{}


    \subsection*{Reitgerte}


    \begin{tabular}{@{}ll@{}}
        \textbf{\trans{weight_label}:} & 1.00 \\
        \textbf{\trans{price_label}:} & 20 \\
        \textbf{\trans{rarity_label}:} & Gewöhnlich \\
        \textbf{\trans{concealment_label}:} & 0 \\
    \end{tabular}

    \wbif{phasesix}{
        \vspace{3mm}
\faIcon{cube}\hspace{0.5em}    }{}


    \subsection*{Eisensporen}


    \begin{tabular}{@{}ll@{}}
        \textbf{\trans{weight_label}:} & 1.00 \\
        \textbf{\trans{price_label}:} & 10 \\
        \textbf{\trans{rarity_label}:} & Gewöhnlich \\
        \textbf{\trans{concealment_label}:} & 0 \\
    \end{tabular}

    \wbif{phasesix}{
        \vspace{3mm}
\faIcon{cube}\hspace{0.5em}    }{}


    \subsection*{Silbersporen}


    \begin{tabular}{@{}ll@{}}
        \textbf{\trans{weight_label}:} & 1.00 \\
        \textbf{\trans{price_label}:} & 50 \\
        \textbf{\trans{rarity_label}:} & Gewöhnlich \\
        \textbf{\trans{concealment_label}:} & 0 \\
    \end{tabular}

    \wbif{phasesix}{
        \vspace{3mm}
\faIcon{cube}\hspace{0.5em}    }{}


    \subsection*{Falknerhandschuh}


    \begin{tabular}{@{}ll@{}}
        \textbf{\trans{weight_label}:} & 2.00 \\
        \textbf{\trans{price_label}:} & 40 \\
        \textbf{\trans{rarity_label}:} & Gewöhnlich \\
        \textbf{\trans{concealment_label}:} & 0 \\
    \end{tabular}

    \wbif{phasesix}{
        \vspace{3mm}
\faIcon{cube}\hspace{0.5em}    }{}


    \subsection*{Maulkorb}


    \begin{tabular}{@{}ll@{}}
        \textbf{\trans{weight_label}:} & 1.00 \\
        \textbf{\trans{price_label}:} & 20 \\
        \textbf{\trans{rarity_label}:} & Gewöhnlich \\
        \textbf{\trans{concealment_label}:} & 0 \\
    \end{tabular}

    \wbif{phasesix}{
        \vspace{3mm}
\faIcon{cube}\hspace{0.5em}    }{}


    \subsection*{Halsband und Leine}

    Halsband und Leine für einen Hund. Oder den Lebensgefährten.

    \begin{tabular}{@{}ll@{}}
        \textbf{\trans{weight_label}:} & 1.00 \\
        \textbf{\trans{price_label}:} & 30 \\
        \textbf{\trans{rarity_label}:} & Gewöhnlich \\
        \textbf{\trans{concealment_label}:} & 0 \\
    \end{tabular}

    \wbif{phasesix}{
        \vspace{3mm}
\faIcon{cube}\hspace{0.5em}    }{}


    \subsection*{Vogelkäfig}


    \begin{tabular}{@{}ll@{}}
        \textbf{\trans{weight_label}:} & 1.00 \\
        \textbf{\trans{price_label}:} & 30 \\
        \textbf{\trans{rarity_label}:} & Gewöhnlich \\
        \textbf{\trans{concealment_label}:} & 5 \\
    \end{tabular}

    \wbif{phasesix}{
        \vspace{3mm}
\faIcon{cube}\hspace{0.5em}    }{}


    \section{Kuriositäten}



    \subsection*{Jade Statue}

    Eine magische Statue eines glatzköpfigen Mannes mit leuchtenden Augen. Bringt man mehrere dieser Statuen zusammen, so wird man in eine seltsame, außerirdisch wirkende Stadt mit hohen Gebäuden und rötlichem Dunst in der Luft teleportert.

    \begin{tabular}{@{}ll@{}}
        \textbf{\trans{weight_label}:} & 1.00 \\
        \textbf{\trans{price_label}:} & 10 \\
        \textbf{\trans{rarity_label}:} & Selten \\
        \textbf{\trans{concealment_label}:} & 0 \\
    \end{tabular}

    \wbif{phasesix}{
        \vspace{3mm}
\faIcon{fire}\hspace{0.5em}    }{}


    \subsection*{Phiole Regenbogenblut von Gargath}

    Eine Phiole voller regenbogenfarbenem Blut von Gargath, dem Wächter des ersten Kreises im verzauberten Wald von Mare.

    \begin{tabular}{@{}ll@{}}
        \textbf{\trans{weight_label}:} & 1.00 \\
        \textbf{\trans{price_label}:} & 1000 \\
        \textbf{\trans{rarity_label}:} & Legendär \\
        \textbf{\trans{concealment_label}:} & 1 \\
        \textbf{\trans{charges_label}:} & 1\\
    \end{tabular}

    \wbif{phasesix}{
        \vspace{3mm}
\faIcon{fire}\hspace{0.5em}    }{}


    \subsection*{Handspiegel}

    Ein einfacher, kleiner Handspiegel

    \begin{tabular}{@{}ll@{}}
        \textbf{\trans{weight_label}:} & 0.30 \\
        \textbf{\trans{price_label}:} & 15 \\
        \textbf{\trans{rarity_label}:} & Gewöhnlich \\
        \textbf{\trans{concealment_label}:} & 0 \\
    \end{tabular}

    \wbif{phasesix}{
        \vspace{3mm}
\faIcon{cube}\hspace{0.5em}    }{}


    \subsection*{Sternenstaub}

    Seltsamer, roter bis regenbogenfarbener Staub einer fremden außerirdischen Stadt. Beim Konsum treten bemerkenswerte Zuwächse an Schnelligkeit und Geschick auf. Nebenwirkungen sind unbekannt.

    \begin{tabular}{@{}ll@{}}
        \textbf{\trans{weight_label}:} & 1.00 \\
        \textbf{\trans{price_label}:} & 10 \\
        \textbf{\trans{rarity_label}:} & Selten \\
        \textbf{\trans{concealment_label}:} & 0 \\
    \end{tabular}

    \wbif{phasesix}{
        \vspace{3mm}
\faIcon{fire}\hspace{0.5em}    }{}


    \subsection*{Kaugummi}


    \begin{tabular}{@{}ll@{}}
        \textbf{\trans{weight_label}:} & 0.01 \\
        \textbf{\trans{price_label}:} & 1 \\
        \textbf{\trans{rarity_label}:} & Gewöhnlich \\
        \textbf{\trans{concealment_label}:} & 0 \\
    \end{tabular}

    \wbif{phasesix}{
        \vspace{3mm}
\faIcon{microchip}\hspace{0.5em}    }{}


    \subsection*{Weihwasserflasche}

    Eine Flasche voller Weihwasser.

Wird es gegen Untote, Vampire oder Werwölfe verwendet, so verursacht es 1W6 Treffer mit Durchschlag 0.

    \begin{tabular}{@{}ll@{}}
        \textbf{\trans{weight_label}:} & 1.00 \\
        \textbf{\trans{price_label}:} & 20 \\
        \textbf{\trans{rarity_label}:} & Gewöhnlich \\
        \textbf{\trans{concealment_label}:} & 0 \\
    \end{tabular}

    \wbif{phasesix}{
        \vspace{3mm}
\faIcon{spider}\hspace{0.5em}    }{}


    \subsection*{Aglaran für Anfänger}

    Dieses futuristische Buch wurde bei einer besonderen Gelegenheit von den Aglaranern erworben. Es vermittelt die aglaraner Sprache für Besucher von fremden Welten.

    \begin{tabular}{@{}ll@{}}
        \textbf{\trans{weight_label}:} & 0.30 \\
        \textbf{\trans{price_label}:} & 1000 \\
        \textbf{\trans{rarity_label}:} & Selten \\
        \textbf{\trans{concealment_label}:} & 0 \\
    \end{tabular}

    \wbif{phasesix}{
        \vspace{3mm}
\faIcon{fire}\hspace{0.5em}    }{}


    \subsection*{Kruzifix}

    Ein christliches Kreuz (optional mit Jesus dran genagelt), welches man in einer Hand halten kann.

Wird es in Sichtweite eines Vampirs gebracht, muss der Vampir einen Resistenz Wurf machen. Misslingt der Wurf, so hat dieser in der folgenden Kampfrunde keine Aktionen.

    \begin{tabular}{@{}ll@{}}
        \textbf{\trans{weight_label}:} & 0.30 \\
        \textbf{\trans{price_label}:} & 10 \\
        \textbf{\trans{rarity_label}:} & Gewöhnlich \\
        \textbf{\trans{concealment_label}:} & 0 \\
    \end{tabular}

    \wbif{phasesix}{
        \vspace{3mm}
\faIcon{spider}\hspace{0.5em}    }{}


    \subsection*{Beauty Set}

    Lippenstift, Kajal und Puder

    \begin{tabular}{@{}ll@{}}
        \textbf{\trans{weight_label}:} & 0.26 \\
        \textbf{\trans{price_label}:} & 19 \\
        \textbf{\trans{rarity_label}:} & Gewöhnlich \\
        \textbf{\trans{concealment_label}:} & 0 \\
    \end{tabular}

    \wbif{phasesix}{
        \vspace{3mm}
\faIcon{microchip}\hspace{0.5em}    }{}


    \subsection*{antike Schivone (Kaufvertrag)}


    \begin{tabular}{@{}ll@{}}
        \textbf{\trans{weight_label}:} & 1.00 \\
        \textbf{\trans{price_label}:} & 1 \\
        \textbf{\trans{rarity_label}:} & Gewöhnlich \\
        \textbf{\trans{concealment_label}:} & 0 \\
    \end{tabular}

    \wbif{phasesix}{
        \vspace{3mm}
\faIcon{fire}\hspace{0.5em}    }{}


    \subsection*{Shoggothenzahn}

    Zahn einer Oma-Shoggothe

    \begin{tabular}{@{}ll@{}}
        \textbf{\trans{weight_label}:} & 1.00 \\
        \textbf{\trans{price_label}:} & 10 \\
        \textbf{\trans{rarity_label}:} & Gewöhnlich \\
        \textbf{\trans{concealment_label}:} & 0 \\
    \end{tabular}

    \wbif{phasesix}{
        \vspace{3mm}
\faIcon{umbrella}\hspace{0.5em}\faIcon{industry}\hspace{0.5em}\faIcon{ankh}\hspace{0.5em}\faIcon{syringe}\hspace{0.5em}\faIcon{spider}\hspace{0.5em}    }{}


    \subsection*{Haarnadel}

    Kann auch als einfacher Dietrich und Stichwerkzeug dienen.

    \begin{tabular}{@{}ll@{}}
        \textbf{\trans{weight_label}:} & 0.03 \\
        \textbf{\trans{price_label}:} & 19 \\
        \textbf{\trans{rarity_label}:} & Gewöhnlich \\
        \textbf{\trans{concealment_label}:} & 0 \\
    \end{tabular}

    \wbif{phasesix}{
        \vspace{3mm}
\faIcon{chess-rook}\hspace{0.5em}\faIcon{umbrella}\hspace{0.5em}\faIcon{flag}\hspace{0.5em}\faIcon{plane}\hspace{0.5em}\faIcon{microchip}\hspace{0.5em}    }{}


    \subsection*{Arlingtonstein}

    Ein Faustgroßer Obsidianstein mit einer flachen Vorderseite, auf der seltsame altertümliche Symbole manchmal glühend zu sehen sind. Werden die Worte in der alten Schrift gesprochen, kann ein altes oder höllisches Wesen beschworen werden.

    \begin{tabular}{@{}ll@{}}
        \textbf{\trans{weight_label}:} & 0.10 \\
        \textbf{\trans{price_label}:} & 1 \\
        \textbf{\trans{rarity_label}:} & Einmalig \\
        \textbf{\trans{concealment_label}:} & 0 \\
    \end{tabular}

    \wbif{phasesix}{
        \vspace{3mm}
\faIcon{spider}\hspace{0.5em}    }{}


    \subsection*{Ouija Board}

    Das Ouija wird von Anhängern des Spiritismus als Hilfsmittel zur Kontaktaufnahme mit Geistern angesehen.

    \begin{tabular}{@{}ll@{}}
        \textbf{\trans{weight_label}:} & 1.00 \\
        \textbf{\trans{price_label}:} & 200 \\
        \textbf{\trans{rarity_label}:} & Außergewöhnlich \\
        \textbf{\trans{concealment_label}:} & 1 \\
    \end{tabular}

    \wbif{phasesix}{
        \vspace{3mm}
\faIcon{spider}\hspace{0.5em}    }{}


    \subsection*{Ring, Silber}

    Ein silberner Ring

    \begin{tabular}{@{}ll@{}}
        \textbf{\trans{weight_label}:} & 0.10 \\
        \textbf{\trans{price_label}:} & 10 \\
        \textbf{\trans{rarity_label}:} & Gewöhnlich \\
        \textbf{\trans{concealment_label}:} & 0 \\
    \end{tabular}

    \wbif{phasesix}{
        \vspace{3mm}
\faIcon{cube}\hspace{0.5em}    }{}


    \subsection*{Stoffpuppe}

    Eine einfache Stoffpuppe.

    \begin{tabular}{@{}ll@{}}
        \textbf{\trans{weight_label}:} & 0.30 \\
        \textbf{\trans{price_label}:} & 10 \\
        \textbf{\trans{rarity_label}:} & Gewöhnlich \\
        \textbf{\trans{concealment_label}:} & 0 \\
    \end{tabular}

    \wbif{phasesix}{
        \vspace{3mm}
\faIcon{cube}\hspace{0.5em}    }{}


    \subsection*{Taschenuhr}


    \begin{tabular}{@{}ll@{}}
        \textbf{\trans{weight_label}:} & 0.08 \\
        \textbf{\trans{price_label}:} & 150 \\
        \textbf{\trans{rarity_label}:} & Gewöhnlich \\
        \textbf{\trans{concealment_label}:} & 0 \\
    \end{tabular}

    \wbif{phasesix}{
        \vspace{3mm}
\faIcon{umbrella}\hspace{0.5em}\faIcon{flag}\hspace{0.5em}\faIcon{plane}\hspace{0.5em}\faIcon{microchip}\hspace{0.5em}    }{}


    \subsection*{Tabakdose}

    Eine Dose um darin Tabak aufzubewaren.

    \begin{tabular}{@{}ll@{}}
        \textbf{\trans{weight_label}:} & 0.30 \\
        \textbf{\trans{price_label}:} & 20 \\
        \textbf{\trans{rarity_label}:} & Gewöhnlich \\
        \textbf{\trans{concealment_label}:} & 0 \\
    \end{tabular}

    \wbif{phasesix}{
        \vspace{3mm}
\faIcon{umbrella}\hspace{0.5em}\faIcon{flag}\hspace{0.5em}\faIcon{plane}\hspace{0.5em}\faIcon{microchip}\hspace{0.5em}    }{}


    \subsection*{Logbuch des Kapitäns}


    \begin{tabular}{@{}ll@{}}
        \textbf{\trans{weight_label}:} & 1.00 \\
        \textbf{\trans{price_label}:} & 10 \\
        \textbf{\trans{rarity_label}:} & Gewöhnlich \\
        \textbf{\trans{concealment_label}:} & 0 \\
    \end{tabular}

    \wbif{phasesix}{
        \vspace{3mm}
\faIcon{dice-six}\hspace{0.5em}\faIcon{umbrella}\hspace{0.5em}\faIcon{car}\hspace{0.5em}\faIcon{flask}\hspace{0.5em}\faIcon{spider}\hspace{0.5em}    }{}


    \subsection*{Haube des logischen Denkens}

    +1 Logik

    \begin{tabular}{@{}ll@{}}
        \textbf{\trans{weight_label}:} & 1.00 \\
        \textbf{\trans{price_label}:} & 1111 \\
        \textbf{\trans{rarity_label}:} & Legendär \\
        \textbf{\trans{concealment_label}:} & 0 \\
    \end{tabular}

    \wbif{phasesix}{
        \vspace{3mm}
\faIcon{cube}\hspace{0.5em}\faIcon{eye}\hspace{0.5em}\faIcon{dice-six}\hspace{0.5em}\faIcon{chess-rook}\hspace{0.5em}\faIcon{umbrella}\hspace{0.5em}\faIcon{flag}\hspace{0.5em}\faIcon{plane}\hspace{0.5em}\faIcon{microchip}\hspace{0.5em}\faIcon{space-shuttle}\hspace{0.5em}\faIcon{industry}\hspace{0.5em}\faIcon{syringe}\hspace{0.5em}\faIcon{9}\hspace{0.5em}\faIcon{book-open}\hspace{0.5em}\faIcon{dragon}\hspace{0.5em}\faIcon{ankh}\hspace{0.5em}\faIcon{magic}\hspace{0.5em}\faIcon{spider}\hspace{0.5em}\faIcon{fire}\hspace{0.5em}    }{}


    \subsection*{Ring, Gold}

    Ein goldener Ring.

    \begin{tabular}{@{}ll@{}}
        \textbf{\trans{weight_label}:} & 0.10 \\
        \textbf{\trans{price_label}:} & 60 \\
        \textbf{\trans{rarity_label}:} & Außergewöhnlich \\
        \textbf{\trans{concealment_label}:} & 0 \\
    \end{tabular}

    \wbif{phasesix}{
        \vspace{3mm}
\faIcon{cube}\hspace{0.5em}    }{}


    \subsection*{Szepter des Lichts}


    \begin{tabular}{@{}ll@{}}
        \textbf{\trans{weight_label}:} & 0.60 \\
        \textbf{\trans{price_label}:} & 1000 \\
        \textbf{\trans{rarity_label}:} & Einmalig \\
        \textbf{\trans{concealment_label}:} & 1 \\
    \end{tabular}

    \wbif{phasesix}{
        \vspace{3mm}
\faIcon{fire}\hspace{0.5em}    }{}


    \subsection*{Sonnenuhr}

    Eine transportable Sonnenuhr.

    \begin{tabular}{@{}ll@{}}
        \textbf{\trans{weight_label}:} & 0.50 \\
        \textbf{\trans{price_label}:} & 20 \\
        \textbf{\trans{rarity_label}:} & Gewöhnlich \\
        \textbf{\trans{concealment_label}:} & 0 \\
    \end{tabular}

    \wbif{phasesix}{
        \vspace{3mm}
\faIcon{cube}\hspace{0.5em}    }{}


    \subsection*{Früchtekuchen}


    \begin{tabular}{@{}ll@{}}
        \textbf{\trans{weight_label}:} & 0.30 \\
        \textbf{\trans{price_label}:} & 10 \\
        \textbf{\trans{rarity_label}:} & Gewöhnlich \\
        \textbf{\trans{concealment_label}:} & 0 \\
    \end{tabular}

    \wbif{phasesix}{
        \vspace{3mm}
\faIcon{cube}\hspace{0.5em}    }{}


    \subsection*{Jonglierbälle}

    Entwerder du kannst es, oder du kannst es nicht.

    \begin{tabular}{@{}ll@{}}
        \textbf{\trans{weight_label}:} & 1.00 \\
        \textbf{\trans{price_label}:} & 10 \\
        \textbf{\trans{rarity_label}:} & Gewöhnlich \\
        \textbf{\trans{concealment_label}:} & 0 \\
    \end{tabular}

    \wbif{phasesix}{
        \vspace{3mm}
\faIcon{cube}\hspace{0.5em}    }{}


    \subsection*{Geisterfalle}

    Ein altes Gerät, das zwei Geister oder ein geisterähnliche Wesen in seinem Umkreis fangen kann. Die Falle muss gestellt werden (Mechanik Wurf) und kann einen Geist fangen. Es gibt einen Mechanismus, um den Geist zu befreien.

    \begin{tabular}{@{}ll@{}}
        \textbf{\trans{weight_label}:} & 1.00 \\
        \textbf{\trans{price_label}:} & 500 \\
        \textbf{\trans{rarity_label}:} & Selten \\
        \textbf{\trans{concealment_label}:} & 2 \\
        \textbf{\trans{charges_label}:} & 2\\
    \end{tabular}

    \wbif{phasesix}{
        \vspace{3mm}
\faIcon{spider}\hspace{0.5em}    }{}


    \subsection*{Brille}

    Eine Brille, hoffentlich auf deine Sehstärke abgestimmt.

    \begin{tabular}{@{}ll@{}}
        \textbf{\trans{weight_label}:} & 0.40 \\
        \textbf{\trans{price_label}:} & 80 \\
        \textbf{\trans{rarity_label}:} & Gewöhnlich \\
        \textbf{\trans{concealment_label}:} & 0 \\
    \end{tabular}

    \wbif{phasesix}{
        \vspace{3mm}
\faIcon{cube}\hspace{0.5em}    }{}


    \subsection*{Teleporter Helm}

    Teleportiert auf das Schiff von KWARG.

    \begin{tabular}{@{}ll@{}}
        \textbf{\trans{weight_label}:} & 1.00 \\
        \textbf{\trans{price_label}:} & 10 \\
        \textbf{\trans{rarity_label}:} & Einmalig \\
        \textbf{\trans{concealment_label}:} & 0 \\
    \end{tabular}

    \wbif{phasesix}{
        \vspace{3mm}
\faIcon{fire}\hspace{0.5em}    }{}


    \subsection*{Goldene Taschenuhr}

    Eine goldene Taschenuhr an einer Kette.

    \begin{tabular}{@{}ll@{}}
        \textbf{\trans{weight_label}:} & 1.00 \\
        \textbf{\trans{price_label}:} & 100 \\
        \textbf{\trans{rarity_label}:} & Gewöhnlich \\
        \textbf{\trans{concealment_label}:} & 0 \\
    \end{tabular}

    \wbif{phasesix}{
        \vspace{3mm}
\faIcon{umbrella}\hspace{0.5em}\faIcon{flag}\hspace{0.5em}\faIcon{plane}\hspace{0.5em}\faIcon{microchip}\hspace{0.5em}    }{}


    \subsection*{Märchenbuch}

    Ein Buch mit Märchen.

    \begin{tabular}{@{}ll@{}}
        \textbf{\trans{weight_label}:} & 1.00 \\
        \textbf{\trans{price_label}:} & 10 \\
        \textbf{\trans{rarity_label}:} & Gewöhnlich \\
        \textbf{\trans{concealment_label}:} & 0 \\
    \end{tabular}

    \wbif{phasesix}{
        \vspace{3mm}
\faIcon{cube}\hspace{0.5em}    }{}


    \subsection*{Historische Bibel}

    Eine gebundene, historische Ausgabe der Bibel.

    \begin{tabular}{@{}ll@{}}
        \textbf{\trans{weight_label}:} & 1.00 \\
        \textbf{\trans{price_label}:} & 100 \\
        \textbf{\trans{rarity_label}:} & Gewöhnlich \\
        \textbf{\trans{concealment_label}:} & 0 \\
    \end{tabular}

    \wbif{phasesix}{
        \vspace{3mm}
\faIcon{chess-rook}\hspace{0.5em}\faIcon{umbrella}\hspace{0.5em}\faIcon{flag}\hspace{0.5em}\faIcon{plane}\hspace{0.5em}\faIcon{microchip}\hspace{0.5em}    }{}


    \subsection*{Goldenes Monokel}

    Ein goldenes Monokel, welches man zum Zweck der guten Sicht vor ein Auge einsetzen kann.

    \begin{tabular}{@{}ll@{}}
        \textbf{\trans{weight_label}:} & 1.00 \\
        \textbf{\trans{price_label}:} & 150 \\
        \textbf{\trans{rarity_label}:} & Gewöhnlich \\
        \textbf{\trans{concealment_label}:} & 0 \\
    \end{tabular}

    \wbif{phasesix}{
        \vspace{3mm}
\faIcon{umbrella}\hspace{0.5em}\faIcon{flag}\hspace{0.5em}\faIcon{dragon}\hspace{0.5em}\faIcon{fire}\hspace{0.5em}    }{}


    \section{Zutaten}



    \subsection*{Schlafmohn (Papaver somniferum)}

    Der Schlafmohn (Papaver somniferum) gehört zur Familie der Mohngewächse und blickt auf eine lange Geschichte als eine unserer ältesten Heilpflanzen zurück.

Er ist jedoch nicht nur medizinisch interessant: Seine Samen werden in der Küche als Nahrungsmittel geschätzt oder zur Gewinnung von hochwertigem Speiseöl gepresst.

Das Besondere am Schlafmohn ist, dass fast alle Teile der Pflanze Alkaloide wie Morphin enthalten. Diese Wirkstoffe konzentrieren sich vor allem im weißen Milchsaft, der durch ein feines Kanalsystem die gesamte Pflanze durchläuft – besonders dicht ist dieses Netz in der Wand der Kapselfrucht. Um an diese Stoffe zu gelangen, werden die unreifen Kapseln leicht eingeritzt, sodass der Saft austreten kann. Wenn dieser Milchsaft an der Luft trocknet, entsteht das bekannte Betäubungsmittel Opium. Der Name leitet sich übrigens aus dem Griechischen ab und bedeutet schlicht „Säftchen“.

    \begin{tabular}{@{}ll@{}}
        \textbf{\trans{weight_label}:} & 10.00 \\
        \textbf{\trans{price_label}:} & 100 \\
        \textbf{\trans{rarity_label}:} & Außergewöhnlich \\
        \textbf{\trans{concealment_label}:} & 0 \\
    \end{tabular}

    \wbif{phasesix}{
        \vspace{3mm}
\faIcon{umbrella}\hspace{0.5em}\faIcon{flag}\hspace{0.5em}\faIcon{plane}\hspace{0.5em}\faIcon{microchip}\hspace{0.5em}\faIcon{spider}\hspace{0.5em}    }{}


    \subsection*{Beifuß (Artemisia vulgaris)}

    Eine Beifuß Pflanze. Die Spitzen des Sprosses werden verwendet um den Gallenfluss und die Verdauung in Gang zu bringen.

    \begin{tabular}{@{}ll@{}}
        \textbf{\trans{weight_label}:} & 0.10 \\
        \textbf{\trans{price_label}:} & 5 \\
        \textbf{\trans{rarity_label}:} & Gewöhnlich \\
        \textbf{\trans{concealment_label}:} & 0 \\
    \end{tabular}

    \wbif{phasesix}{
        \vspace{3mm}
\faIcon{cube}\hspace{0.5em}\faIcon{flask}\hspace{0.5em}    }{}


    \subsection*{Eibisch (Althaea officinalis)}

    Von dieser Heilpflanze wird die Wurzel verwendet. Diese wird kalt angesetzt und muss ungefähr zwei Stunden ziehen. Erst nach dem Ziehen wird die Flüssigkeit gesiebt und anschließend erhitzt. Die Stoffe bieten Schutz für die Schleimhäute und wirken reizlindernd. Eine hilfreiche Heilpflanze bei Magen-Darm-Problemen und Husten.

    \begin{tabular}{@{}ll@{}}
        \textbf{\trans{weight_label}:} & 0.10 \\
        \textbf{\trans{price_label}:} & 3 \\
        \textbf{\trans{rarity_label}:} & Gewöhnlich \\
        \textbf{\trans{concealment_label}:} & 0 \\
    \end{tabular}

    \wbif{phasesix}{
        \vspace{3mm}
\faIcon{cube}\hspace{0.5em}\faIcon{flask}\hspace{0.5em}    }{}


    \subsection*{Kamille (Matricaria recutita)}

    Kamille zählt zu den ältesten Heilpflanzen und wurde bereits im Mittelalter angewandt. Die Blüten haben eine heilende und beruhigende Wirkung. Äußerlich kann Kamille bei Entzündungen des Zahnfleisches, der Haut oder der Schleimhaut angewendet werden. Innerlich eingenommen wirkt sie bei Magen-Darm-Erkrankungen. Spülen und die Inhalation sind ebenfalls weit verbreitet.

    \begin{tabular}{@{}ll@{}}
        \textbf{\trans{weight_label}:} & 0.10 \\
        \textbf{\trans{price_label}:} & 2 \\
        \textbf{\trans{rarity_label}:} & Gewöhnlich \\
        \textbf{\trans{concealment_label}:} & 0 \\
    \end{tabular}

    \wbif{phasesix}{
        \vspace{3mm}
\faIcon{cube}\hspace{0.5em}\faIcon{flask}\hspace{0.5em}    }{}


    \subsection*{Salbei (Salvia officinalis)}

    Die Blätter des Salbeis haben eine entzündungshemmende, schweißhemmende und zusammenziehende Wirkung. Ein Tee oder Spülungen sind bei Halsentzündungen oder auch Schweißausbrüchen empfehlenswert.

    \begin{tabular}{@{}ll@{}}
        \textbf{\trans{weight_label}:} & 0.10 \\
        \textbf{\trans{price_label}:} & 5 \\
        \textbf{\trans{rarity_label}:} & Gewöhnlich \\
        \textbf{\trans{concealment_label}:} & 0 \\
    \end{tabular}

    \wbif{phasesix}{
        \vspace{3mm}
\faIcon{cube}\hspace{0.5em}\faIcon{flask}\hspace{0.5em}    }{}


    \subsection*{Schmetterlingsdrachensekret}

    Geht man vorsichtig zu Werke, so lassen sich Schmetterlingsdrachen melken. Sie sondern ein gar seltsames Sekret ab, welches denjenigen, der es konsumiert sofort in einen Schlaf mit faszinierenden Träumen fallen lässt.

Wird der Trank verabreicht oder genommen, so schläft die konsumierende Person mindestens acht Stunden tief und fest. Für diese Zeit wird die doppelte Rast angewendet. Der Schlafende ist allenfalls durch echte Schmerzen zu wecken.

    \begin{tabular}{@{}ll@{}}
        \textbf{\trans{weight_label}:} & 0.10 \\
        \textbf{\trans{price_label}:} & 200 \\
        \textbf{\trans{rarity_label}:} & Selten \\
        \textbf{\trans{concealment_label}:} & 0 \\
    \end{tabular}

    \wbif{phasesix}{
        \vspace{3mm}
\faIcon{dragon}\hspace{0.5em}\faIcon{flask}\hspace{0.5em}\faIcon{spider}\hspace{0.5em}    }{}


    \subsection*{Ammoniak}

    Ammoniak (NH\_3) fungiert als essenzieller Basischmikalie in der chemischen Industrie.

Das Hauptanwendungsgebiet liegt in der Produktion von stickstoffhaltigen Düngemitteln (z. B. Harnstoff oder Ammoniumnitrat), welche den Großteil des weltweiten Verbrauchs ausmachen. Darüber hinaus dient Ammoniak als Vorläufer für organische Synthesen, insbesondere bei der Herstellung von Polymeren und synthetischen Fasern.

Die Substanz ist als toxisch eingestuft. In hoher Konzentration wirken Ammoniakdämpfe stark reizend auf die Okularschleimhäute sowie das respiratorische System, wobei schwere Lungenschäden bis hin zum letalen Ausgang möglich sind. In der täglichen Anwendung sind solche Extremereignisse jedoch selten.

    \begin{tabular}{@{}ll@{}}
        \textbf{\trans{weight_label}:} & 20.00 \\
        \textbf{\trans{price_label}:} & 50 \\
        \textbf{\trans{rarity_label}:} & Gewöhnlich \\
        \textbf{\trans{concealment_label}:} & 0 \\
    \end{tabular}

    \wbif{phasesix}{
        \vspace{3mm}
\faIcon{umbrella}\hspace{0.5em}\faIcon{flag}\hspace{0.5em}\faIcon{plane}\hspace{0.5em}\faIcon{microchip}\hspace{0.5em}\faIcon{space-shuttle}\hspace{0.5em}\faIcon{spider}\hspace{0.5em}    }{}


    \subsection*{Kräutermischung}

    Eine leckere Kräutermischung um Essen zu verfeinern.

    \begin{tabular}{@{}ll@{}}
        \textbf{\trans{weight_label}:} & 0.10 \\
        \textbf{\trans{price_label}:} & 5 \\
        \textbf{\trans{rarity_label}:} & Gewöhnlich \\
        \textbf{\trans{concealment_label}:} & 0 \\
        \textbf{\trans{charges_label}:} & 10\\
    \end{tabular}

    \wbif{phasesix}{
        \vspace{3mm}
\faIcon{cube}\hspace{0.5em}\faIcon{flask}\hspace{0.5em}    }{}


    \subsection*{Bernstein}

    Ein glatter, ovaler Bernstein mit einem warmen, goldenen Schimmer. Seine polierte Oberfläche ist leicht transparent und reflektiert das Licht auf faszinierende Weise. Der handgroße Stein wirkt durch seine geschwungene Form wie ein natürlicher Talisman.

    \begin{tabular}{@{}ll@{}}
        \textbf{\trans{weight_label}:} & 0.10 \\
        \textbf{\trans{price_label}:} & 50 \\
        \textbf{\trans{rarity_label}:} & Außergewöhnlich \\
        \textbf{\trans{concealment_label}:} & 0 \\
    \end{tabular}

    \wbif{phasesix}{
        \vspace{3mm}
\faIcon{cube}\hspace{0.5em}    }{}


    \subsection*{Brennnessel (Urtica dioica)}

    Brennnesseln haben eine entwässernde und entzündungshemmende Wirkung. Ein Tee aus den Blättern der Brennnessel verschafft Linderung bei Rheuma und Gicht.

    \begin{tabular}{@{}ll@{}}
        \textbf{\trans{weight_label}:} & 0.10 \\
        \textbf{\trans{price_label}:} & 2 \\
        \textbf{\trans{rarity_label}:} & Gewöhnlich \\
        \textbf{\trans{concealment_label}:} & 0 \\
    \end{tabular}

    \wbif{phasesix}{
        \vspace{3mm}
\faIcon{cube}\hspace{0.5em}    }{}


    \subsection*{Schlüsselblume (Primula veris)}

    Die Schlüsselblume war im als Fruchtbarkeits- und Schutzmittel bekannt. Heute hilft ein Schlüsselwurzeltee gegen Erkältungen. Salbei und Fenchel verstärken die Wirkung.

    \begin{tabular}{@{}ll@{}}
        \textbf{\trans{weight_label}:} & 0.10 \\
        \textbf{\trans{price_label}:} & 5 \\
        \textbf{\trans{rarity_label}:} & Gewöhnlich \\
        \textbf{\trans{concealment_label}:} & 0 \\
    \end{tabular}

    \wbif{phasesix}{
        \vspace{3mm}
\faIcon{cube}\hspace{0.5em}    }{}


    \subsection*{Schafgarbe (Achillea millefolium)}

    Schafgarbe wird wegen ihrer blutstillenden Wirkung eingesetzt. Die Blüten und die Blätter enthalten Gerb-, Bitter- und Mineralstoffe. Das ätherische Öl der Pflanze wirkt entzündungshemmend und krampflösend.

    \begin{tabular}{@{}ll@{}}
        \textbf{\trans{weight_label}:} & 0.10 \\
        \textbf{\trans{price_label}:} & 3 \\
        \textbf{\trans{rarity_label}:} & Gewöhnlich \\
        \textbf{\trans{concealment_label}:} & 0 \\
    \end{tabular}

    \wbif{phasesix}{
        \vspace{3mm}
\faIcon{cube}\hspace{0.5em}    }{}


    \subsection*{Calciumchlorid}

    Bei Calciumchlorid handelt es sich um eine Verbindung aus Calcium und Chlor mit der Formel CaCl\_2.

Chemisch gesehen befindet sich das Calcium darin in der Oxidationsstufe +2 und das Chlor in der Stufe -1. Man findet diesen Stoff in der Natur vor allem in Salzsolen gelöst. Wenn Calciumchlorid Wasser bindet und fest wird, entstehen daraus sehr seltene Minerale: Das Dihydrat wird als Sinjarit bezeichnet, während das Hexahydrat unter dem Namen Antarcticit bekannt ist. Man kann es jedoch auch einfach online bestellen.

    \begin{tabular}{@{}ll@{}}
        \textbf{\trans{weight_label}:} & 200.00 \\
        \textbf{\trans{price_label}:} & 50 \\
        \textbf{\trans{rarity_label}:} & Gewöhnlich \\
        \textbf{\trans{concealment_label}:} & 0 \\
    \end{tabular}

    \wbif{phasesix}{
        \vspace{3mm}
\faIcon{umbrella}\hspace{0.5em}\faIcon{flag}\hspace{0.5em}\faIcon{plane}\hspace{0.5em}\faIcon{microchip}\hspace{0.5em}\faIcon{spider}\hspace{0.5em}    }{}


    \subsection*{Lavendel (Lavandula officinalis)}

    Im elften Jahrhundert wurde Lavendel von Mönchen in Mitteleuropa angesiedelt. In der Medizin wurde Lavendel eine Wirksamkeit bei Insektenstichen und Verbrennungen zugeschrieben. Ein Lavendel Tee hilft bei Erkältungen und Kopfschmerzen.

    \begin{tabular}{@{}ll@{}}
        \textbf{\trans{weight_label}:} & 0.10 \\
        \textbf{\trans{price_label}:} & 4 \\
        \textbf{\trans{rarity_label}:} & Gewöhnlich \\
        \textbf{\trans{concealment_label}:} & 0 \\
    \end{tabular}

    \wbif{phasesix}{
        \vspace{3mm}
\faIcon{cube}\hspace{0.5em}\faIcon{flask}\hspace{0.5em}    }{}


    \subsection*{Arnika (Arnica montana)}

    Arnika wird bei Entzündungen, Wunden, um den Kreislauf anzuregen und als Abtreibungsmittel verwendet. Die Blüten werden als Salbe, als Tee oder als Tinktur verwendet.

    \begin{tabular}{@{}ll@{}}
        \textbf{\trans{weight_label}:} & 0.10 \\
        \textbf{\trans{price_label}:} & 5 \\
        \textbf{\trans{rarity_label}:} & Gewöhnlich \\
        \textbf{\trans{concealment_label}:} & 0 \\
    \end{tabular}

    \wbif{phasesix}{
        \vspace{3mm}
\faIcon{cube}\hspace{0.5em}    }{}


    \subsection*{Goldnugget}

    Ein kleines Stück unverarbeitetes Gold, etwa 5 Gramm.

    \begin{tabular}{@{}ll@{}}
        \textbf{\trans{weight_label}:} & 0.05 \\
        \textbf{\trans{price_label}:} & 300 \\
        \textbf{\trans{rarity_label}:} & Gewöhnlich \\
        \textbf{\trans{concealment_label}:} & 0 \\
    \end{tabular}

    \wbif{phasesix}{
        \vspace{3mm}
\faIcon{cube}\hspace{0.5em}    }{}


    \subsection*{Melisse (Melissa officinalis)}

    Melisse wurde seit jeher als Heilkraut in der Medizin genutzt. Sie wirkt gegen Kopfschmerzen, Nervosität, Schlafstörungen und Magen-Darm-Beschwerden. Zudem bringt ein Aufguss mit Melisse Entspannung.

    \begin{tabular}{@{}ll@{}}
        \textbf{\trans{weight_label}:} & 0.10 \\
        \textbf{\trans{price_label}:} & 2 \\
        \textbf{\trans{rarity_label}:} & Gewöhnlich \\
        \textbf{\trans{concealment_label}:} & 0 \\
    \end{tabular}

    \wbif{phasesix}{
        \vspace{3mm}
\faIcon{cube}\hspace{0.5em}\faIcon{flask}\hspace{0.5em}    }{}


    \subsection*{Engelwurzen (Angelica archangelica)}

    Die Pflanze wird bei Verdauungsschwäche, Appetitlosigkeit und Verdauungsschwäche angewendet, und soll angeblich vor der Pest schützen.

    \begin{tabular}{@{}ll@{}}
        \textbf{\trans{weight_label}:} & 0.10 \\
        \textbf{\trans{price_label}:} & 3 \\
        \textbf{\trans{rarity_label}:} & Gewöhnlich \\
        \textbf{\trans{concealment_label}:} & 0 \\
    \end{tabular}

    \wbif{phasesix}{
        \vspace{3mm}
\faIcon{cube}\hspace{0.5em}\faIcon{flask}\hspace{0.5em}    }{}


    \subsection*{Alant (Inula helenium)}

    Diese Heilpflanze aus dem Mittelalter ist in modernen Zeiten nicht mehr weit verbreitet. Ihre Anwendung verbessert die Verdauung, und ihr wird eine vorbeugende Wirkung gegen Darmkrebs zugeschrieben.

    \begin{tabular}{@{}ll@{}}
        \textbf{\trans{weight_label}:} & 0.10 \\
        \textbf{\trans{price_label}:} & 5 \\
        \textbf{\trans{rarity_label}:} & Gewöhnlich \\
        \textbf{\trans{concealment_label}:} & 0 \\
    \end{tabular}

    \wbif{phasesix}{
        \vspace{3mm}
\faIcon{cube}\hspace{0.5em}\faIcon{flask}\hspace{0.5em}    }{}


    \subsection*{Worse than life}

    DIE Droge der Zukunft. Das Crystal Meth der Zukunft. Weiss, stark und unwiderstehlich.

    \begin{tabular}{@{}ll@{}}
        \textbf{\trans{weight_label}:} & 0.10 \\
        \textbf{\trans{price_label}:} & 5 \\
        \textbf{\trans{rarity_label}:} & Außergewöhnlich \\
        \textbf{\trans{concealment_label}:} & 0 \\
    \end{tabular}

    \wbif{phasesix}{
        \vspace{3mm}
\faIcon{space-shuttle}\hspace{0.5em}\faIcon{fire}\hspace{0.5em}    }{}


    \subsection*{Beinwell (Symphytum officinale)}

    Beinwell regt die Durchblutung an, Prellungen, Hämatome und Verstauchungen verschwinden schneller. Beinwell beschleunigt die Regeneration der Zellen.

    \begin{tabular}{@{}ll@{}}
        \textbf{\trans{weight_label}:} & 0.10 \\
        \textbf{\trans{price_label}:} & 5 \\
        \textbf{\trans{rarity_label}:} & Gewöhnlich \\
        \textbf{\trans{concealment_label}:} & 0 \\
    \end{tabular}

    \wbif{phasesix}{
        \vspace{3mm}
\faIcon{cube}\hspace{0.5em}    }{}


    \subsection*{Thymian (Thymus vulgaris)}

    Thymian wird seit über 4000 Jahren gegen Keuchhusten, Husten und Bronchitis eingesetzt. Seine schleimlösende Wirkung wird besonders geschätzt.

    \begin{tabular}{@{}ll@{}}
        \textbf{\trans{weight_label}:} & 0.10 \\
        \textbf{\trans{price_label}:} & 5 \\
        \textbf{\trans{rarity_label}:} & Gewöhnlich \\
        \textbf{\trans{concealment_label}:} & 0 \\
    \end{tabular}

    \wbif{phasesix}{
        \vspace{3mm}
\faIcon{cube}\hspace{0.5em}\faIcon{flask}\hspace{0.5em}    }{}


    \subsection*{Schöllkraut (Chelidonium majus)}

    Im Mittelalter wurde Schöllkraut bei Hautausschlägen, Sehschwäche oder Gelbsucht verwendet. Die Alkaloide der Pflanze habenden einen krampflösenden Effekt. Sie helfen bei Verdauungsproblemen und regen den Gallenfluss an.

    \begin{tabular}{@{}ll@{}}
        \textbf{\trans{weight_label}:} & 0.10 \\
        \textbf{\trans{price_label}:} & 5 \\
        \textbf{\trans{rarity_label}:} & Gewöhnlich \\
        \textbf{\trans{concealment_label}:} & 0 \\
    \end{tabular}

    \wbif{phasesix}{
        \vspace{3mm}
\faIcon{cube}\hspace{0.5em}\faIcon{flask}\hspace{0.5em}    }{}


    \subsection*{Baldrian (Valeriana officinalis)}

    Baldrian hilft bei Einschlafstörungen und Unruhe. Hopfen und Melisse steigern die Wirkung des Baldrians und verbessern den Geschmack.

    \begin{tabular}{@{}ll@{}}
        \textbf{\trans{weight_label}:} & 0.10 \\
        \textbf{\trans{price_label}:} & 3 \\
        \textbf{\trans{rarity_label}:} & Gewöhnlich \\
        \textbf{\trans{concealment_label}:} & 0 \\
    \end{tabular}

    \wbif{phasesix}{
        \vspace{3mm}
\faIcon{cube}\hspace{0.5em}\faIcon{flask}\hspace{0.5em}    }{}


    \subsection*{Spitzwegerich (Plantago lanceolata)}

    Die spitzen, schmalen Blätter des Spitzwegerich werden als Sirup oder auch als Tee bei Erkältungskrankheiten verwendet. Spitzwegerich kann auch zerkleinert und zerrieben auf Wunden oder Insektenstiche aufgetragen werden und entfaltet hier eine kühlende Wirkung. Die Pflanze wird auch bei Durchfall eingesetzt.

    \begin{tabular}{@{}ll@{}}
        \textbf{\trans{weight_label}:} & 0.10 \\
        \textbf{\trans{price_label}:} & 2 \\
        \textbf{\trans{rarity_label}:} & Gewöhnlich \\
        \textbf{\trans{concealment_label}:} & 0 \\
    \end{tabular}

    \wbif{phasesix}{
        \vspace{3mm}
\faIcon{cube}\hspace{0.5em}\faIcon{flask}\hspace{0.5em}    }{}


    \section{Musikinstrumente}



    \subsection*{Laute}

    Die Laute (über spätmittelhochdeutsch lūte von arabisch العود, DMG al-ʿūd ‚der Stab, das Holz, Laute‘) ist ein Zupfinstrument mit Korpus und angesetztem Hals sowie mit gleichlaufend zur Instrumentendecke verlaufenden Saiten.

    \begin{tabular}{@{}ll@{}}
        \textbf{\trans{weight_label}:} & 1.50 \\
        \textbf{\trans{price_label}:} & 25 \\
        \textbf{\trans{rarity_label}:} & Gewöhnlich \\
        \textbf{\trans{concealment_label}:} & 3 \\
    \end{tabular}

    \wbif{phasesix}{
        \vspace{3mm}
\faIcon{eye}\hspace{0.5em}\faIcon{chess-rook}\hspace{0.5em}\faIcon{umbrella}\hspace{0.5em}\faIcon{flag}\hspace{0.5em}\faIcon{plane}\hspace{0.5em}\faIcon{microchip}\hspace{0.5em}\faIcon{ankh}\hspace{0.5em}\faIcon{9}\hspace{0.5em}    }{}


    \subsection*{Dudelsack}

    Zwergischer Dudelsack. Da ist genug Luft im Behälter um zu Singen während Zwerg spielt.

    \begin{tabular}{@{}ll@{}}
        \textbf{\trans{weight_label}:} & 5.00 \\
        \textbf{\trans{price_label}:} & 250 \\
        \textbf{\trans{rarity_label}:} & Gewöhnlich \\
        \textbf{\trans{concealment_label}:} & 0 \\
    \end{tabular}

    \wbif{phasesix}{
        \vspace{3mm}
\faIcon{chess-rook}\hspace{0.5em}\faIcon{umbrella}\hspace{0.5em}\faIcon{flag}\hspace{0.5em}\faIcon{plane}\hspace{0.5em}\faIcon{microchip}\hspace{0.5em}\faIcon{dragon}\hspace{0.5em}\faIcon{magic}\hspace{0.5em}\faIcon{spider}\hspace{0.5em}\faIcon{fire}\hspace{0.5em}    }{}

\end{multicols}